% ---------------------------------------------------------------------------
% Chapter 4: Cousins Nobody Invited: Underlying Event and Colour Reconnection
% The scripts, records and figures discussed here live in example/MPI/ of
% https://github.com/ravindkv/JitJet . The book gives the physics and the
% instructions; the code itself is read on GitHub, not printed here.
% ---------------------------------------------------------------------------
%% chapter-local shorthands (\pt already carries a subscript, so it cannot take another)
\providecommand{\ptzero}{\ensuremath{p_{\mathrm{T}0}}\xspace}
\providecommand{\sigint}{\ensuremath{\sigma_{\mathrm{int}}}\xspace}
\providecommand{\signd}{\ensuremath{\sigma_{\mathrm{ND}}}\xspace}
\providecommand{\bmean}{\ensuremath{\langle b\rangle}\xspace}
\chapter{Cousins Nobody Invited: Underlying Event and Colour Reconnection}
\chaptermark{Cousins Nobody Invited}
\label{ch:uninvited_cousins}

\begin{journeybox}
The rest of the two protons does not stay at home. Other partons scatter softly,
beam remnants hadronise, and colour strings reconnect across the event. Distant
cousins join the road and some of their energy ends up inside our jet.
\end{journeybox}

\physicsline{the divergent QCD $2\to2$ cross section and its regularisation,
multiparton interactions as a Poisson process in impact parameter, the
overlap of two protons and the enhancement factor, the chain of interactions
below the hard scale, depleted parton densities and companion quarks, beam
remnants and primordial transverse momentum, the MPI-based colour
reconnection and the string length, the transverse-region definition of the
underlying event, the CP5 tune, the underlying-event energy inside a cone}

\section{Welcome back: what we will see at this third stop}
\label{sec:uninvited_cousins:what_we_compute}

At the end of \cref{ch:mother_gives_birth} the record contained forty
partons and a hole. The two incoming partons of \file{event_PS.lhe} carried
momentum fractions $x_A=0.017$ and $x_B=0.47$ and nothing balanced them;
the two protons that had supplied them were simply absent. That is not how
a collision ends. Each proton is a bag of partons, and when two bags pass
through each other the pair that produced our $\bbbar$ is only the hardest
of the pairs that scatter. The others scatter too, at lower transverse
momentum, each with its own small shower, and what is left of each proton,
the \emph{beam remnant}, flies on down the beam pipe carrying most of the
energy of the collision. Between all of them run colour strings, and the
strings do not respect which scattering a parton came from. This stop is
about the family that was already on the road when ours set out: the
cousins nobody invited.

The number to take away is small and it is the first entry in the ledger
that \emph{adds} to the mother's family rather than subtracting from it. In
our event, no parton from the rest of the protons happened to land within
$\dR<0.4$ of the $b$ quark, so the ledger row of the previous chapter
survives to the hundredth of a GeV; within $\dR<0.8$ the underlying event
adds $3.4\GeV$ and lifts the family mass from $13.5$ to $20.2\GeV$. On
average, over many dressings of the same hard event at the same collision
geometry, the $R=0.4$ cone receives $1.6\pm2.2\GeV$ and the $R=0.8$ cone
$5.9\pm4.4\GeV$. The uninvited cousins are a few percent of the jet at
$R=0.4$ and are the reason jets are not built with $R=1$; the detector will
later add two hundred \emph{other} collisions on top
(\cref{ch:crowd_on_the_road}), and the methods of
\cref{ch:paying_for_uninvited} will have to tell the two apart. As at the
previous stops we meet the number four times, from the top down:
\begin{enumerate}[nosep]
  \item \textbf{Watch it.} A film shows two protons overlapping, the chain
    of secondary scatterings running down in \pt, the remnants being formed
    from what is left, the strings reconnecting, and the same hard event
    dressed four hundred times (\cref{sec:uninvited_cousins:animation}).
  \item \textbf{Run it.} A Python file, needing only NumPy and SciPy, takes
    the showered record of \cref{ch:mother_gives_birth} and adds the
    multiparton interactions, the beam remnants and the colour
    reconnection, in two seconds, following \PYTHIA's construction class by
    class (\cref{sec:uninvited_cousins:standalone_code}).
  \item \textbf{Derive it.} Why the QCD $2\to2$ cross section must exceed
    the proton--proton cross section, how that paradox is resolved by
    counting scatterings per collision, how the impact parameter enters,
    how the chain is generated with the veto algorithm of
    \cref{sec:mother_gives_birth:veto}, what the protons look like after each
    scattering, how the remnants are formed, and what colour reconnection
    does to the strings
    (\crefrange{sec:uninvited_cousins:the_rest_of_the_protons_does_not_stay_qu}{sec:uninvited_cousins:measuring_the_underlying_event_transvers}).
  \item \textbf{Check it against the official packages.} A validation
    script dresses the same hard event two thousand times with the real
    \PYTHIA, and the CMS generator job of \cref{ch:ancestral_home} already
    had all of this switched on, so its log and its third figure are read
    with the new status codes in hand
    (\cref{sec:uninvited_cousins:validation}).
\end{enumerate}
\MG again does not appear: multiparton interactions, remnants and colour
reconnection are \PYTHIA's job in every CMS sample, whatever produced the
hard process.

\Cref{tab:uninvited_cousins:numbers} gives away the ending, with the same
warning as \cref{tab:mother_gives_birth:numbers} and one more. The stage is
random, so single events are compared on structure and ensembles on
averages. The new warning is the \emph{impact parameter}: the first thing
the model draws is how centrally the two protons hit, and nothing in the
hard process fixes it. Seed~1 of the standalone script, left to itself,
draws a peripheral collision with a single secondary scattering, which would
make a dull chapter; the \CMSSW event of \cref{ch:ancestral_home} happened to
be central, $b=0.31\bmean$, with 31 secondary scatterings. The book's record
therefore fixes the impact parameter to the \CMSSW value, so that the two
can be compared at the same geometry, and every ensemble below is given both
ways: with $b$ drawn, as nature does it, and at $b=0.3136\bmean$.

\begin{table}[htbp]\centering\small
\caption[The rest of the protons obtained three ways]{The rest of the
protons obtained three ways. ``Secondary'' means not belonging to the hard
$\bbbar$ system and its shower; ``UE'' (underlying event) means the partons
of the secondary systems and the remnants; cones are drawn around the final
copy of the $b$ quark. The first four columns are averages over $2000$
(\PYTHIA), $400$ (standalone), $537$ (\PYTHIA at fixed $b$) and $400$
dressings of the same showered event, the spread of a single event after
$\pm$; the fifth is the one dressing with seed~1 at $b=0.3136\bmean$ that
wrote \file{example/MPI/event_MPI.lhe}; the last is the single \CMSSW event.
The standalone ensembles keep the seed-1 shower of
\cref{ch:mother_gives_birth} and redraw only this stage; \PYTHIA redraws
everything.}
\label{tab:uninvited_cousins:numbers}
\scriptsize\setlength{\tabcolsep}{3.5pt}
\begin{tabular}{lrrrrrr}\toprule
 & \PYTHIA & standalone & \PYTHIA & standalone & standalone & \CMSSW \\
 & $b$ drawn & $b$ drawn & $b=0.3136$ & $b=0.3136$ & seed 1 & 1 event \\ \midrule
interactions, incl.\ the hard one & $18.0\pm9.9$ & $14.7\pm7.2$ & $27.5\pm5.4$ & $21.4\pm4.2$ & $25$ & $32$ \\
$b/\bmean$, $f(b)$ & $0.61$, $2.26$ & $0.63$, $2.10$ & $0.31$, $3.70$ & $0.31$, $3.63$ & $0.31$, $3.63$ & $0.31$, $3.71$ \\
final partons & $125\pm60$ & $138\pm54$ & $181\pm34$ & $190\pm33$ & $211$ & $169$ \\
\quad of which from secondary systems & $92$ & $88$ & $145$ & $137$ & $155$ & -- \\
\quad of which beam remnants & $10.2$ & $9.4$ & $13.6$ & $12.9$ & $16$ & $13$ \\
$\sum\pt$ of the secondary partons [GeV] & $134$ & $124$ & $212$ & $196$ & $230$ & -- \\
hardest secondary interaction [GeV] & $6.8$ & $6.5$ & $8.6$ & $8.4$ & $11.45$ & $6.65$ \\ \midrule
transverse UE density [GeV per unit $\eta$--$\phi$] & $2.21\pm1.72$ & $2.00\pm1.59$ & $3.48\pm1.54$ & $3.23\pm1.44$ & $2.10$ & -- \\
UE $\sum\pt$ within $\dR<0.4$ of the $b$ [GeV] & $1.21\pm2.48$ & $0.85\pm1.31$ & $1.80\pm2.43$ & $1.60\pm2.19$ & $0.00$ & -- \\
UE $\sum\pt$ within $\dR<0.8$ of the $b$ [GeV] & $4.57\pm5.20$ & $3.95\pm4.05$ & $7.32\pm5.68$ & $5.88\pm4.44$ & $3.42$ & -- \\
UE partons within $\dR<0.4$ / $0.8$ & $0.75$ / $2.99$ & $0.72$ / $2.87$ & $1.20$ / $4.75$ & $1.07$ / $4.18$ & $0$ / $3$ & -- \\ \midrule
systems merged by colour reconnection & -- & $13.2$ & -- & $19.7$ & $23$ & -- \\
string length $\lambda$ before $\to$ after & -- & $517\to326$ & -- & $708\to406$ & $806\to403$ & -- \\
\bottomrule
\end{tabular}
\end{table}

\paragraph{A picture to carry: the other tenants.} The project picture of
\cref{ch:mother_gives_birth} had two main contractors and their
subcontractors. This chapter adds the rest of the building. The hard process
rented the top floor; the other floors are rented too, by tenants who never
heard of our project, each running a small project of their own with the
same rules, the same subcontracting, the same recoils. They share the
building's budget (the protons' momentum), so every tenant that moves in
leaves less for the next, and the landlord (the remnant) keeps what nobody
rented. Then the electricians come: the colour strings that hadronisation
will pull into hadrons do not run floor by floor but are rewired so that
the total length of cable is as short as possible, and a cable that starts
on our floor may now end on somebody else's. The picture captures the three
things the chapter derives: more than one scattering per collision, a
shared and depleted budget, and strings that cross system boundaries. Its
limit is the same as before: the scatterings are independent only in the
model's leading approximation, and the model has knobs, the tune, which
the data fix.

\section{Watch the protons empty out: the animated journey}
\label{sec:uninvited_cousins:animation}

The film first. \file{example/MPI/Standalone/animate_standalone_multiparton_interactions.py}
imports the standalone script of \cref{sec:uninvited_cousins:standalone_code}
and replays the seed-1 dressing at the \CMSSW impact parameter in five scenes
of about fifty seconds in all, once without and once with colour
reconnection, so that every number on the screen is produced by the classes
that wrote the record. It is on the JitJet playlist,
\begin{center}
  \url{https://www.youtube.com/playlist?list=PLFtvDH9g5Km0}
\end{center}
and \cref{fig:uninvited_cousins:stills} shows one key frame per scene. Each
scene names the section that explains it.
\begin{enumerate}
  \item \textbf{Two protons meet} (\cref{sec:uninvited_cousins:impact}).
    Two discs of matter, the CP5 double-Gaussian profile with $63\%$ of the
    matter in a core of $0.76$ of the radius, slide towards each other in
    the transverse plane. As the impact parameter $b$ shrinks the overlap
    $O(b)$ and the enhancement $f(b)=n(b)/\langle n\rangle$ on the right
    rise along their curves; the collision stops at $b=0.3136\bmean$, the
    value of the \CMSSW event, where $f=3.63$ and the Poisson mean of
    secondary scatterings, $32.5$, appears. Watch how fast $f$ falls: at
    $b=\bmean$ it is already $0.52$, and the seed-1 draw without the
    \code{--impact} option, $b=1.21\bmean$, sits at $f=0.20$.
  \item \textbf{The chain runs down in \pt} (\cref{sec:uninvited_cousins:chain}).
    A scale runs from the hard scale, $113\GeV$, to $\pt^{\min}=0.2\GeV$.
    Every time it passes the \pt of one of the 24 secondary scatterings the
    scattering happens: the ladder on the left gains a rung, the
    $(\eta,\phi)$ picture in the middle gains a pair of cyan squares among
    the partons of the hard system, and the no-interaction curve on the
    right shows how unlikely the quiet stretch from $113$ down to
    $11.5\GeV$ was. The rungs crowd together below $2\GeV$: seventeen of
    the 24 are there.
  \item \textbf{What is left} (\cref{sec:uninvited_cousins:remnants}). The
    momentum bars of the two protons are eaten interaction by interaction,
    red for the hard process, cyan for the secondary systems; the valence
    content changes when a valence quark leaves, and at the end the grey
    remnant partons appear with their flavours: a $ud_0$ diquark and six
    companion quarks for proton~A, a $ud_1$ and eight companions for
    proton~B. On the right the primordial-$k_{\mathrm{T}}$ kicks of every
    initiator and remnant are drawn as arrows; per proton they sum to zero.
  \item \textbf{The strings reconnect} (\cref{sec:uninvited_cousins:colour_reconnection_and_its_effect_on_je}).
    The colour strings of the event in the $(y,\phi)$ plane, before and
    after the MPI-based colour reconnection. Before, each of the 25 systems
    has its own strings and they criss-cross the whole rapidity range;
    after, the gluons of 23 soft systems have been moved onto the strings
    of harder systems and the total string length $\lambda$ has dropped
    from $806$ to $403$. The $b$ quark's own string is highlighted: it is
    one of the 51 that remain.
  \item \textbf{Dressed again and again} (\cref{tab:uninvited_cousins:numbers}).
    The same showered event dressed with seeds $1,2,3,\ldots$ at the same
    impact parameter: the histograms of the number of interactions, of the
    transverse underlying-event density and of the extra \pt inside
    $\dR<0.4$ around the $b$ fill up, and their running means settle on
    $21.4$, $3.23\GeV$ and $1.60\GeV$ next to the \PYTHIA values $27.5$,
    $3.48$ and $1.80$ of the validation run. Our event, the first frame, is
    on the lucky side for the cone and on the busy side for the chain.
\end{enumerate}
Rendering needs matplotlib in addition to NumPy and SciPy, and
\code{ffmpeg} for an \code{.mp4}:
\begin{minted}[linenos=false,frame=single,framesep=1.5mm,fontsize=\small]{bash}
cd example/MPI/Standalone
python3 animate_standalone_multiparton_interactions.py            # the film, about a minute to render
python3 animate_standalone_multiparton_interactions.py --fast     # low-resolution preview
python3 animate_standalone_multiparton_interactions.py --scene chain -o chain.gif
python3 animate_standalone_multiparton_interactions.py --stills stills.pdf --no-video
\end{minted}
The physics options are those of the standalone script: \code{--seed 7}
films another dressing, \code{--impact 1.0} a peripheral collision, and
\code{--no-cr} leaves the strings as they were.

\begin{figure}[htbp]\centering
\includegraphics[page=1,width=0.49\textwidth]{MPI/Standalone/animate_standalone_multiparton_interactions_stills.pdf}\hfill
\includegraphics[page=2,width=0.49\textwidth]{MPI/Standalone/animate_standalone_multiparton_interactions_stills.pdf}\\[2mm]
\includegraphics[page=3,width=0.49\textwidth]{MPI/Standalone/animate_standalone_multiparton_interactions_stills.pdf}\hfill
\includegraphics[page=4,width=0.49\textwidth]{MPI/Standalone/animate_standalone_multiparton_interactions_stills.pdf}\\[2mm]
\includegraphics[page=5,width=0.49\textwidth]{MPI/Standalone/animate_standalone_multiparton_interactions_stills.pdf}
\caption[Key frames of the animated standalone MPI stage]{Key frames of the
five scenes of the animated standalone MPI stage, drawn by
\file{animate_standalone_multiparton_interactions.py} with \code{--stills}.
Top left: the two protons at impact parameter $b$ with the overlap and the
enhancement factor. Top right: the \pt ladder of the 24 secondary
scatterings, the $(\eta,\phi)$ picture and the no-interaction probability.
Middle left: the momentum bars of the two protons being used up and the
primordial-$k_{\mathrm{T}}$ kicks. Middle right: the colour strings before
and after reconnection and the string length. Bottom: four hundred
dressings of the same event against the \PYTHIA reference. The film is on
the JitJet YouTube playlist.}
\label{fig:uninvited_cousins:stills}
\end{figure}

\section{Press the button: the standalone MPI stage}
\label{sec:uninvited_cousins:standalone_code}

Now dress the event yourself. The material of this chapter is in
\file{example/MPI/} of the JitJet repository,
\begin{center}
  \url{https://github.com/ravindkv/JitJet}
\end{center}
with the script, its log, its validation and its plots under
\file{example/MPI/Standalone/}, the event record and the plotting script one
level up, and the \CMSSW figure under \file{example/MPI/CMSSW/}. As before
the code is not printed; each script gets a description, the command and
the output to expect.

\subsection{The standalone script}

The script, \file{standalone_multiparton_interactions.py} in
\file{example/MPI/Standalone/}, is a port of four classes of \PYTHIA~8,
\code{MultipartonInteractions},
\code{BeamParticle}, \code{BeamRemnants} and the \code{reconnectMPIs} model of
\code{ColourReconnection}, into two thousand lines of Python that need NumPy
and SciPy and nothing else. It imports the shower script of
\cref{ch:mother_gives_birth} for the running coupling, the event record and
the shower itself, which it runs on every secondary system; its own PDF
grid, embedded between \code{BEGIN/END GENERATED TABLES} by
\file{dump_mpi_tables.py}, is the same NNPDF3.1 set on LHAPDF's knots but
reaches down to $x=10^{-9}$, because a $0.2\GeV$ scattering at rapidity~9
takes $x\sim10^{-8}$ from a proton. A map of the file: \code{SigmaInt} (the
regularised cross section, \cref{sec:uninvited_cousins:cross_section});
\code{Overlap} (the proton profile, $O(b)$ and $f(b)$,
\cref{sec:uninvited_cousins:impact}); \code{Beam} and \code{BeamPDF} (the
depleted proton, \cref{sec:uninvited_cousins:flavours}); \code{Generator}
with \code{next\_pt}, \code{trial} and \code{scatter} (the chain,
\cref{sec:uninvited_cousins:chain}), \code{shower\_system} (the chapter-2
shower with the depleted proton on each side), \code{remnants}
(\cref{sec:uninvited_cousins:remnants}), \code{reconnect},
\code{remnant\_colours} and \code{lambda\_measure}
(\cref{sec:uninvited_cousins:colour_reconnection_and_its_effect_on_je});
and \code{summarise}, which measures the underlying event and the cones
around the two $b$ quarks. From the repository root:
\begin{minted}[linenos=false,frame=single,framesep=1.5mm,fontsize=\small]{bash}
cd example/MPI/Standalone
python3 standalone_multiparton_interactions.py --impact 0.3136      # seed 1 -> ../event_MPI.lhe (the book's record)
python3 standalone_multiparton_interactions.py                      # draw the impact parameter instead
python3 standalone_multiparton_interactions.py --impact 0.3136 -v   # learning mode: every number, every step
python3 standalone_multiparton_interactions.py --impact 0.3136 -vv  # ... every veto trial and the secondary showers
python3 standalone_multiparton_interactions.py --no-cr --no-kt      # colour reconnection and primordial kT off
python3 standalone_multiparton_interactions.py --repeat 400 --impact 0.3136 --quiet --history s.json
\end{minted}
The first command integrates the cross section at start-up (three seconds)
and generates one event (two seconds, most of it the 24 secondary showers),
then prints
\begin{minted}[linenos=false,frame=single,framesep=1.5mm,fontsize=\footnotesize]{text}
pT0 = 1.4417 GeV, sigma_int = 497.33 mb, sigma_ND = 55.42 mb, <n> = 8.974; f(b) ranges up to 4.52, <f>_hard = 2.218
input example/PS/event_PS.lhe: 963 rows after MPI; seed = 1
impact parameter b = 0.3136 <b> with enhancement f(b) = 3.627; expected 32.5 interactions, got 25 (1 hard + 24 MPI) ...
  pT of the secondary interactions [GeV]: 11.45 7.68 4.47 3.65 3.42 3.14 2.80 2.19 2.18 1.79 1.76 1.63 1.50 ...
  final state: 211 partons, of which 155 from the MPI systems (incl. their showers: 31 ISR + 76 FSR branchings) ...
  underlying event (|eta| < 2.5, w.r.t. the leading b): toward 49.47, transverse 21.99, away 47.67 GeV; ...
  bbar: pT = 72.50 GeV at eta = -3.42; cone R = 0.4: 4 partons, pT 91.27 (hard system) -> 93.88 GeV with 1 UE ...
  b: pT = 78.29 GeV at eta = -1.38; cone R = 0.4: 3 partons, pT 100.67 (hard system) -> 100.67 GeV with 0 UE ...
  colour reconnection: 23 systems merged; string length lambda = 805.7 -> 403.5; 51 colour-singlet strings
\end{minted}
and writes the record to \file{example/MPI/event_MPI.lhe}. With
\code{--history} the chain and the summary go to a JSON file as well; with
\code{--repeat N} the script dresses the showered event $N$ times with
seeds $1\ldots N$ and writes the per-event quantities of
\cref{tab:uninvited_cousins:numbers} to \file{s.stats.json}, about fifteen
minutes for $400$. The physics options are \code{--pt0ref} and
\code{--ecmpow} (the regulator of \cref{sec:uninvited_cousins:cross_section}
and its energy dependence), \code{--sigma-nd}, \code{--impact} (the
impact parameter in units of $\bmean$, drawn when absent), \code{--cr-range}
(the reconnection strength), \code{--ptmin-isr}, \code{--ptmin-fsr} and
\code{--pt0} (the cut-offs of the secondary showers), \code{--no-shower},
\code{--no-isr}, \code{--no-fsr}, \code{--no-cr}, \code{--no-kt} and
\code{--radius}. The output with \code{--seed 1 --impact 0.3136} is
byte-for-byte the committed record.

The two learning modes work as in \cref{ch:mother_gives_birth}. \code{-v}
prints, before the event, the regularised cross section at a dozen scales
with its integral, the overlap table in $b$, and the state of the two
protons after the hard process; then, for every secondary interaction, the
accepted scale and its veto weight, the rapidities and momentum fractions,
the depleted densities of both protons at that $x$, the six largest channel
weights, the chosen process with its colour flow and momenta, and the size
of its shower; then the remnant flavours, the primordial-$k_{\mathrm{T}}$
widths and kicks, the light-cone solution, every remnant row, every colour
reconnection decision with the dipole each gluon was inserted into, the
string length before and after, the conservation checks and the full
record. \code{-vv} adds every trial scale of the veto algorithm and the
secondary showers' own \code{-v} output. The $1500$ lines of \code{-v} for
seed~1 at $b=0.3136\bmean$ are committed as
\file{example/MPI/Standalone/standalone_multiparton_interactions.log}, and
every number quoted in the theory sections is taken from that file. The
plots of the theory sections come from a sibling script,
\begin{minted}[linenos=false,frame=single,framesep=1.5mm,fontsize=\small]{bash}
python3 example/MPI/Standalone/plot_theory_MPI.py   # -> theory_multiparton_interactions.pdf, one page per section
\end{minted}
which replays the seed-1 dressing with the classes of the script and reads
the \code{--repeat} statistics and the \PYTHIA validation files when they
are present; \code{make figures} regenerates it.

\subsection{Reading the record}
\label{sec:uninvited_cousins:record_format}

The record keeps \PYTHIA's conventions and adds five status codes to those
of \cref{sec:mother_gives_birth:record_format}. A secondary scattering is
written like the hard process: $-31$ for its two incoming partons and $-33$
for its two outgoing ones, negative because every one of them is replaced by
the shower's copies. At the end of the stage every system is boosted to the
frame of its initiators after the primordial $k_{\mathrm{T}}$ and the
remnant kinematics; the boosted initiators are written with status $-61$ and
the boosted final partons with status~$62$, so that the final state of the
record is every row with status~$62$ plus the beam remnants, status~$63$.
Our record has $963$ rows, of which $195$ are status~$62$ and $16$ are
status~$63$; the mother of a status-$62$ row is its last copy of the shower
stage, so the family of any final parton is still found by walking up the
mother column, now through a $-31$ or $-33$ row if the parton belongs to a
secondary system. The momentum sum at the bottom of the file is, for the
first time, that of the two protons: $(E,p_x,p_y,p_z)=(13600,0,0,0)\GeV$ to
two MeV.

\subsection{Drawing the event}

\file{example/MPI/plot_journey_MPI.py} draws the record in the $(\eta,\phi)$
plane in the layout of the \CMSSW analyser of
\cref{sec:uninvited_cousins:cmssw} (\cref{fig:uninvited_cousins:journey}),
reusing the reader and the frame of chapter~2's script and adding the two
origins of this stage, secondary scatterings and beam remnants. It applies
the analyser's classification rule on purpose: a copy or an FSR radiator
inherits the origin of its mother, so a gluon \emph{emitted} by a parton of
a secondary system is counted as FSR and an initial-state emission of a
secondary system as ISR, and only the outgoing partons of the scatterings
themselves and their descendants by recoil are counted as ``MPI''. That is
why the figure says $45$ MPI partons while the script's summary says $155$
partons from the secondary systems: the two numbers count the same partons
with different rules, and \cref{sec:uninvited_cousins:cmssw} needs the
analyser's.
\begin{minted}[linenos=false,frame=single,framesep=1.5mm,fontsize=\small]{bash}
python3 example/MPI/plot_journey_MPI.py                 # -> example/MPI/journey_MPI.pdf
python3 example/MPI/plot_journey_MPI.py -i my.lhe --png
\end{minted}

\section{The rest of the protons does not stay quiet: multiparton interactions}
\label{sec:uninvited_cousins:the_rest_of_the_protons_does_not_stay_qu}

\subsection{A cross section larger than the proton: the regularised \texorpdfstring{$2\to2$}{2 to 2} rate}
\label{sec:uninvited_cousins:cross_section}

Guess first. \Cref{ch:ancestral_home} computed the cross section for
$\pthat>100\GeV$ of our process, $345\pb$. What is the cross section for
\emph{any} QCD $2\to2$ scattering with $\pthat>2\GeV$? The matrix elements
of \cref{sec:ancestral_home:cmssw} fall as $1/\pthat^4$ and the parton
densities rise steeply at small $x$, so the answer is enormous: with the
numbers of this chapter, $168\,\mathrm{mb}$, that is $1.7\times10^{11}\pb$,
and above $0.2\GeV$ it is $497\,\mathrm{mb}$. The total proton--proton cross
section at $13.6\TeV$ is $100\,\mathrm{mb}$, and the part of it that is not
elastic or diffractive, the non-diffractive cross section \signd, is
$55\,\mathrm{mb}$. A parton-level cross section ten times larger than the
cross section for the two protons to interact at all is not a
contradiction; it is a count. If each non-diffractive collision contains on
average $\langle n\rangle$ parton--parton scatterings above the cut, then
\begin{equation}
  \langle n\rangle=\frac{\sigint(\pt>\pt^{\min})}{\signd},
  \label{eq:uninvited_cousins:nmean}
\end{equation}
and the parton-level number, correctly read, says that an average collision
at the LHC contains \emph{nine} scatterings above $0.2\GeV$, not one. This
is the observation of Sj\"ostrand and van~Zijl~\cite{Sjostrand:1987su} on
which every \PYTHIA event rests, and it has a corollary that the rest of the
section makes precise: our hard scattering is accompanied by the others,
which do not know about it.

\begin{figure}[htbp]\centering
\resizebox{\textwidth}{!}{% ---------------------------------------------------------------------------
% Why there is more than one scattering per collision. Left: the two protons
% of our event pass through each other; the hard u ubar -> b bbar scattering
% is one of many parton-parton scatterings, the others soft. Right: the
% parton-parton cross section dsigma/dpT^2 against pT, bare (1/pT^4, divergent)
% and regularised (1/(pT^2+pT0^2)^2); its integral above pTmin, sigma_int,
% exceeds the non-diffractive cross section, so <n> = sigma_int/sigma_ND
% scatterings happen per collision. Numbers from the -v log of
% standalone_multiparton_interactions.py (section "setting the scene").
% ---------------------------------------------------------------------------
\begin{tikzpicture}[x=1cm,y=1cm,
  every node/.style={font=\footnotesize},
  proton/.style={draw,very thick,fill=gray!10},
  parton/.style={circle,inner sep=0pt,minimum size=2mm},
  hard/.style={line width=1.8pt,red!80!black},
  soft/.style={line width=0.9pt,teal!80!black},
  lab/.style={font=\scriptsize,align=center},
  head/.style={font=\small\bfseries},
  gluon/.style={thick,decorate,decoration={coil,aspect=0,segment length=3pt,amplitude=1.2pt},gray!70}]
% ---------------- (a) one collision, many scatterings ------------------------------------------
\node[head,anchor=west] at (-3.4,3.2) {(a) one collision, 25 parton--parton scatterings};
\node[proton,ellipse,minimum width=46mm,minimum height=30mm] (A) at (-1.4,0) {};
\node[proton,ellipse,minimum width=46mm,minimum height=30mm,fill=gray!14] (B) at (1.4,0.6) {};
\node[lab,anchor=east] at (-3.9,0) {proton A\\$+z$};
\node[lab,anchor=west] at (3.9,0.6) {proton B\\$-z$};
\draw[-{Stealth[length=2mm]},gray!70,thick] (-4.6,-1.6) -- (-3.4,-1.6) node[midway,below,font=\scriptsize] {$6800\GeV$};
\draw[-{Stealth[length=2mm]},gray!70,thick] (4.6,2.2) -- (3.4,2.2) node[midway,above,font=\scriptsize] {$6800\GeV$};
% the hard scattering: u ubar -> b bbar
\coordinate (H) at (0.0,0.3);
\draw[hard] (-2.4,-0.3) -- (H) node[pos=0.35,below,font=\scriptsize,text=red!80!black] {$\bar u$, $x=0.0022$};
\draw[hard] (2.4,0.9) -- (H) node[pos=0.35,above,font=\scriptsize,text=red!80!black] {$u$, $x=0.39$};
\draw[hard,-{Stealth[length=2.2mm]}] (H) -- (1.6,-1.7) node[right,font=\scriptsize,text=red!80!black] {$b$, $113\GeV$};
\draw[hard,-{Stealth[length=2.2mm]}] (H) -- (-1.6,2.3) node[left,font=\scriptsize,text=red!80!black] {$\bar b$};
\node[circle,fill=red!80!black,inner sep=1.6pt] at (H) {};
% the soft scatterings (24 of them): drawn as small crosses with short outgoing legs
\foreach \px/\py/\ang in {-1.1/0.9/20, -0.6/-0.7/-35, 0.7/1.3/60, 0.9/-0.4/-70, -0.2/1.6/15, 0.3/-1.1/40, 1.5/0.2/-20, -1.6/0.1/70, -0.5/0.5/-55, 1.1/0.9/35, 0.0/-0.3/80, -1.0/-1.1/10}
  {\node[circle,fill=teal!80!black,inner sep=1.1pt] at (\px,\py) {};
   \draw[soft,-{Stealth[length=1.4mm]}] (\px,\py) -- ++(\ang:0.55);
   \draw[soft,-{Stealth[length=1.4mm]}] (\px,\py) -- ++(\ang+180:0.55);}
\node[lab,text=teal!80!black,anchor=north] at (0.0,-2.0) {24 secondary scatterings, $\pt$ from $11.5$ down to $0.43\GeV$,\\each with its own shower, drawn from the depleted protons};
\node[lab,text=black!65,anchor=north] at (0.0,-2.8) {what is left of each proton becomes the beam remnants};
% ---------------- (b) the cross section and its regulator --------------------------------------
\begin{scope}[xshift=9.0cm,yshift=-2.6cm]
\node[head,anchor=west] at (-0.3,5.8) {(b) the cross section must be tamed};
\draw[-{Stealth}] (0,0) -- (6.4,0) node[below left=2pt and -2pt] {$\pt$};
\draw[-{Stealth}] (0,0) -- (0,5.2) node[above,align=center,font=\scriptsize] {$\mathrm{d}\sigma/\mathrm{d}\pt^2$};
\foreach \x/\l in {1.5/$p_{\mathrm{T}0}$,3.0/$2p_{\mathrm{T}0}$,4.5/$3p_{\mathrm{T}0}$} {\draw (\x,0) -- (\x,-0.1) node[below] {\l};}
\draw[dashed,red!70!black] (1.5,0) -- (1.5,4.6);
\node[font=\scriptsize,text=red!70!black,anchor=west] at (1.55,4.75) {$p_{\mathrm{T}0}=1.44\GeV$ (CP5 at $13.6\TeV$)};
% bare: 4.5/x^4 in units of pT0 -> x = pT/pT0, plotted from where it enters the frame
\begin{scope}
\clip (0.05,0) rectangle (6.3,4.6);
\draw[very thick,dashed,gray!70!black] plot[domain=0.95:4.2,samples=120,variable=\t] ({1.5*\t},{4.5/(\t*\t*\t*\t)});
\draw[very thick,teal!80!black] plot[domain=0.05:4.2,samples=160,variable=\t] ({1.5*\t},{4.5/((\t*\t+1)*(\t*\t+1))});
\fill[teal!80!black,opacity=0.12] plot[domain=0.13:4.2,samples=120,variable=\t] ({1.5*\t},{4.5/((\t*\t+1)*(\t*\t+1))}) -- (6.3,0) -- (0.2,0) -- cycle;
\end{scope}
\node[font=\scriptsize,text=gray!70!black,anchor=west,align=left] at (2.3,3.9) {bare: $\alphas^2(\pt^2)/\pt^4$\\diverges, $\sigma_{\mathrm{int}}\to\infty$};
\node[font=\scriptsize,text=teal!80!black,anchor=west,align=left] at (2.3,2.6) {regularised: $\dfrac{\alphas^2(\pt^2+p_{\mathrm{T}0}^2)}{(\pt^2+p_{\mathrm{T}0}^2)^2}$\\colour screening below $p_{\mathrm{T}0}$};
\draw[dotted,black!60] (0.2,0) -- (0.2,4.4);
\node[font=\scriptsize,anchor=south,rotate=90] at (0.08,2.0) {$\pt^{\min}=0.2\GeV$};
\node[font=\scriptsize,text=teal!80!black,anchor=north west,align=left] at (0.35,1.5) {shaded: $\sigma_{\mathrm{int}}=497\,\mathrm{mb}$\\$\sigma_{\mathrm{ND}}=55.4\,\mathrm{mb}$\\$\langle n\rangle=\sigma_{\mathrm{int}}/\sigma_{\mathrm{ND}}=9.0$};
\end{scope}
\end{tikzpicture}
}
\caption[One collision, many scatterings, and the regularised cross
section]{Left: one proton--proton collision contains many parton--parton
scatterings; in our event the hard $\bar uu\to\bbbar$ at $113\GeV$ and 24
softer ones between $11.5$ and $0.43\GeV$. Right: the parton--parton cross
section $\mathrm{d}\sigma/\mathrm{d}\pt^2$ diverges as $1/\pt^4$ (dashed)
and is regularised by the replacement of \cref{eq:uninvited_cousins:regularised}
(solid); its integral above $\pt^{\min}$, shaded, is nine times the
non-diffractive cross section.}
\label{fig:uninvited_cousins:two_protons_sketch}
\end{figure}

The count cannot be left to diverge, though, and the $1/\pt^4$ of the
matrix element does diverge as $\pt\to0$. Physically the divergence is cut
off by the finite size of the proton: a gluon of wavelength longer than the
distance between colour charges does not resolve them and is screened. The
model implements the screening by a single scale \ptzero,
\begin{equation}
  \frac{\mathrm{d}\hat\sigma}{\mathrm{d}\pt^2}\propto
  \frac{\alphas^2(\pt^2)}{\pt^4}
  \quad\longrightarrow\quad
  \frac{\alphas^2(\pt^2+\ptzero^2)}{(\pt^2+\ptzero^2)^2},
  \qquad
  \ptzero(E_{\mathrm{cm}})=\ptzero^{\mathrm{ref}}\left(\frac{E_{\mathrm{cm}}}{7\TeV}\right)^{\epsilon},
  \label{eq:uninvited_cousins:regularised}
\end{equation}
applied to every $2\to2$ channel of \cref{ch:ancestral_home} plus the
heavy-flavour channels $gg,q\bar q\to c\bar c,\bbbar$ with massive
kinematics. The regulator grows slowly with energy, because the proton is
resolved at smaller $x$ where the gluon density is higher and the screening
stronger; the two numbers $\ptzero^{\mathrm{ref}}$ and $\epsilon$ are the
most important parameters of any tune, and CP5~\cite{CMS:2019csb} sets them
to $1.41\GeV$ and $0.03344$. The full differential cross section summed over
channels and integrated over the rapidities of the two outgoing partons is
\begin{equation}
  \frac{\mathrm{d}\sigma}{\mathrm{d}\pt^2}=\sum_{i,j}\int\mathrm{d}y_3\,\mathrm{d}y_4\;
  x_1f_i(x_1,Q^2)\,x_2f_j(x_2,Q^2)\;\frac{\mathrm{d}\hat\sigma_{ij}}{\mathrm{d}\hat t}
  \Big|_{\mathrm{regularised}},\qquad
  x_{1,2}=\frac{\pt}{E_{\mathrm{cm}}}\left(e^{\pm y_3}+e^{\pm y_4}\right),
  \label{eq:uninvited_cousins:dsigma}
\end{equation}
with $Q^2=\pt^2+\ptzero^2$ in the densities too, and
$\sigint=\int_{\pt^{\min}}\mathrm{d}\sigma$. Two details decide the number
and both cost me a day: a final state of two different partons is reached
twice when both rapidities run over their full range, once as the
$t$-channel and once as the $u$-channel configuration, so the sum must be
halved for those channels, exactly as \PYTHIA averages the two; and the
densities below the grid's lowest scale, $Q=1.65\GeV$, are frozen rather
than extrapolated, as \PYTHIA's own readers do, which matters because
scatterings below $1.5\GeV$ carry two thirds of \sigint.

\paragraph{The numbers.} The \code{-v} log opens with
$\ptzero=1.41\times(13600/7000)^{0.03344}=1.4417\GeV$ and
$\signd=100.309-21.89-23.0=55.42\,\mathrm{mb}$, the \CMSSW settings for the
total, elastic and diffractive cross sections at $13.6\TeV$, and then the
table of \cref{eq:uninvited_cousins:dsigma}:
\begin{center}\footnotesize
\begin{tabular}{rrrrr}\toprule
$\pt$ [GeV] & $\alphas(\pt^2+\ptzero^2)$ & $\bigl(\pt^2/(\pt^2+\ptzero^2)\bigr)^2$ & $\mathrm{d}\sigma/\mathrm{d}\pt^2$ [mb/GeV$^2$] & $\sigma(>\pt)$ [mb] \\ \midrule
$0.2$ & $0.370$ & $0.0004$ & $599$ & $497.3$ \\
$1.0$ & $0.330$ & $0.106$ & $100$ & $285.6$ \\
$1.442$ & $0.305$ & $0.250$ & $38.7$ & $217.9$ \\
$2.0$ & $0.278$ & $0.433$ & $20.9$ & $168.1$ \\
$5.0$ & $0.211$ & $0.852$ & $1.61$ & $33.95$ \\
$10.0$ & $0.178$ & $0.960$ & $0.0772$ & $5.17$ \\
$20.0$ & $0.153$ & $0.990$ & $0.00236$ & $0.561$ \\
$113.3$ & $0.114$ & $0.9997$ & $1.2\times10^{-7}$ & $0.00075$ \\
\bottomrule
\end{tabular}
\end{center}
The third column is the regulator's effect relative to the bare rate: at
$\pt=\ptzero$ one quarter, at $0.2\GeV$ four parts in ten thousand. The
integral, $\sigint=497.33\,\mathrm{mb}$, is what the \CMSSW log prints as
``\code{pT0 = 1.44 gives sigmaInteraction = 500.67 mb: accepted}'' at
initialisation, and $\langle n\rangle=497.33/55.42=8.97$. The cross section
above our hard scale, $0.75\,\mathrm{\mu b}$, is the last entry: a
$2\to2$ scattering above $113\GeV$ happens in one collision in $70\,000$,
which is why the hard process is always the first rung of the ladder and not
somewhere in the middle. The composition of the rate (panel~(c) of
\cref{fig:uninvited_cousins:cross_section_plot}) is gluon--gluon everywhere
below $10\GeV$: at the scale of our first secondary scattering the log
lists $gg\to gg$ with $65\%$ of the weight and the five next channels, all
quark--gluon, with $2$--$4\%$ each.

\begin{figure}[htbp]\centering
\includegraphics[page=1,width=\textwidth]{MPI/Standalone/theory_multiparton_interactions.pdf}
\caption[The regularised cross section and \sigint]{Page~1 of
\file{theory_multiparton_interactions.pdf}. (a) The bare and the regularised
parton--parton cross section of \cref{eq:uninvited_cousins:regularised}
against \pt, with \ptzero and $\pt^{\min}$ marked. (b) \sigint against the
regulator \ptzero, computed with the script's \code{SigmaInt} on a coarser
grid, the CP5 point against the \CMSSW log, and \signd as the horizontal
line; the right axis is $\langle n\rangle$. (c) The share of the rate from
$gg$, $qg$ and quark--quark initial states against \pt: gluons dominate
below $10\GeV$, where all the secondary scatterings of our event are.}
\label{fig:uninvited_cousins:cross_section_plot}
\end{figure}

\subsection{Where the two protons meet: the impact parameter}
\label{sec:uninvited_cousins:impact}

Guess again. Our hard scattering happened. Does that tell us anything about
how centrally the two protons hit? It does, and the direction is not
obvious: a hard scattering is more likely when more partons overlap, so the
collisions that contain one are more central than average, and a central
collision contains more of everything else as well. The number of
scatterings is therefore not Poisson with mean $\langle n\rangle$ in every
collision; it is Poisson with a mean that depends on the impact parameter,
$n(b)$, and the hard process biases $b$ towards zero. This is the second
ingredient of the model~\cite{Sjostrand:1987su}, and the one that makes the
spread in \cref{tab:uninvited_cousins:numbers} as large as it is.

\begin{figure}[htbp]\centering
\resizebox{\textwidth}{!}{% ---------------------------------------------------------------------------
% The impact parameter. Left: the two protons in the transverse plane, each a
% double-Gaussian matter distribution (CP5: 63 % of the matter in a core of
% 0.76 of the radius), displaced by b; the overlap O(b) is the hatched region.
% Right: the enhancement factor f(b) = n(b)/<n> of the -v log ("the impact
% parameter" section) with our event and the averages marked.
% ---------------------------------------------------------------------------
\begin{tikzpicture}[x=1cm,y=1cm,
  every node/.style={font=\footnotesize},
  lab/.style={font=\scriptsize,align=center},
  head/.style={font=\small\bfseries}]
% ---------------- (a) two protons in the transverse plane ----------------------------------
\node[head,anchor=west] at (-3.6,3.3) {(a) two protons meet at impact parameter $b$};
\coordinate (CA) at (-1.0,0);
\coordinate (CB) at (0.7,0.45);
% halo + core of each proton (double Gaussian): drawn as concentric discs
\foreach \c/\col in {CA/blue,CB/orange} {
  \fill[\col!15] (\c) circle (2.2);
  \fill[\col!30] (\c) circle (1.45);
  \fill[\col!45] (\c) circle (0.85);}
\begin{scope}
\clip (CA) circle (1.45);
\fill[pattern=north east lines,pattern color=black!70] (CB) circle (1.45);
\end{scope}
\draw[very thick,blue!70!black] (CA) circle (2.2);
\draw[very thick,orange!80!black] (CB) circle (2.2);
\node[circle,fill=blue!70!black,inner sep=1.5pt] at (CA) {};
\node[circle,fill=orange!80!black,inner sep=1.5pt] at (CB) {};
\draw[{Stealth}-{Stealth},thick] (CA) -- (CB) node[midway,below right=-1pt,font=\scriptsize] {$b$};
\node[lab,text=blue!70!black] at (-2.6,-2.0) {proton A\\core $63\%$, radius $0.76\,a$};
\node[lab,text=orange!80!black] at (2.4,2.3) {proton B\\(the same)};
\node[lab] at (-0.15,0.2) {};
\node[lab,anchor=north] at (-0.2,-2.6) {hatched: the overlap $O(b)=\int\rho_A(\vec r)\,\rho_B(\vec r-\vec b)\,\mathrm{d}^2r$\\the number of scatterings is Poisson with mean $n(b)=k\,O(b)$};
% ---------------- (b) f(b) --------------------------------------------------------------------
\begin{scope}[xshift=7.4cm,yshift=-2.7cm]
\node[head,anchor=west] at (-0.3,6.0) {(b) the enhancement $f(b)=n(b)/\langle n\rangle$};
\draw[-{Stealth}] (0,0) -- (6.6,0) node[right] {$b/\langle b\rangle$};
\draw[-{Stealth}] (0,0) -- (0,5.4) node[above,font=\scriptsize] {$f(b)$};
\foreach \x/\l in {0/0,1.5/0.5,3/1,4.5/1.5,6/2} {\draw (\x,0) -- (\x,-0.1) node[below] {\l};}
\foreach \y/\l in {1.1/1,2.2/2,3.3/3,4.4/4} {\draw (0,\y) -- (-0.1,\y) node[left] {\l};}
% f(b) from the log: b/<b> 0 .25 .5 .75 1 1.25 1.5 2 -> 4.517 3.928 2.594 1.315 0.522 0.166 0.044 0.002
\draw[very thick,teal!80!black] plot[smooth] coordinates {(0,4.969) (0.75,4.321) (1.5,2.853) (2.25,1.447) (3.0,0.574) (3.75,0.183) (4.5,0.048) (6.0,0.002)};
\draw[dashed,red!80!black] (0.941,0) -- (0.941,3.99);
\node[circle,fill=red!80!black,inner sep=2pt] at (0.941,3.99) {};
\node[font=\scriptsize,text=red!80!black,anchor=west,align=left] at (1.05,4.3) {our event: $b=0.3136\,\langle b\rangle$\\$f=3.63$ (\CMSSW: $3.71$), $n(b)=32.5$};
\draw[dotted,black!60] (0,2.44) -- (6.4,2.44) node[pos=0.98,above left,font=\scriptsize,text=black!65] {$\langle f\rangle_{\mathrm{hard}}=2.22$};
\draw[dotted,black!60] (0,1.07) -- (6.4,1.07) node[pos=0.45,above left,font=\scriptsize,text=black!65] {$\langle f\rangle_{\mathrm{ND}}=0.97$};
\node[font=\scriptsize,text=teal!80!black,anchor=south west] at (0.1,4.97) {$f(0)=4.52$};
\node[font=\scriptsize,anchor=north west,align=left,text=black!65,text width=3.0cm] at (3.3,2.05) {minimum bias weights $b$ by $P_{\mathrm{int}}(b)$; a hard process by $O(b)$: central collisions favoured};
\end{scope}
\end{tikzpicture}
}
\caption[Two protons at impact parameter $b$ and the enhancement
$f(b)$]{Left: the two protons of our event in the transverse plane, each a
double-Gaussian matter distribution, displaced by the impact parameter $b$;
the overlap $O(b)$ is hatched. Right: the enhancement factor $f(b)$ of
\cref{eq:uninvited_cousins:fb} computed by the script, with our event at
$b=0.3136\bmean$, the average over hard events and the average over
minimum-bias events marked.}
\label{fig:uninvited_cousins:overlap_sketch}
\end{figure}

The proton is given a matter distribution; CP5 uses a double Gaussian with
$63\%$ of the matter in a core of radius $0.7634$ of the outer radius
(\code{coreFraction}, \code{coreRadius} in the \CMSSW settings table). The
overlap of two such protons displaced by $\vec b$ is the time-integrated
product of the densities,
\begin{equation}
  O(b)=\int\mathrm{d}t\,\mathrm{d}^3x\;\rho(\vec x)\,\rho(\vec x+\vec b+\vec v t)
  =\sum_{\text{core/halo pairs}}c_{ij}\,\frac{e^{-b^2/(a_i^2+a_j^2)}}{\pi(a_i^2+a_j^2)},
  \label{eq:uninvited_cousins:overlap}
\end{equation}
a sum of three Gaussians in $b$ for the double-Gaussian profile. The mean
number of scatterings at impact parameter $b$ is taken proportional to the
overlap, $n(b)=k\,O(b)$, the probability that the collision interacts at all
is $P_{\mathrm{int}}(b)=1-e^{-n(b)}$, and $k$ is fixed by requiring that the
average over interacting collisions reproduce \cref{eq:uninvited_cousins:nmean},
\begin{equation}
  \langle n\rangle=\frac{\int\mathrm{d}^2b\;n(b)}{\int\mathrm{d}^2b\;P_{\mathrm{int}}(b)}
  =\frac{\sigint}{\signd},
  \qquad
  f(b)\equiv\frac{n(b)}{\langle n\rangle}
  =\frac{O(b)\int\mathrm{d}^2b'\,P_{\mathrm{int}}(b')}{\int\mathrm{d}^2b'\,O(b')}.
  \label{eq:uninvited_cousins:fb}
\end{equation}
The denominator of the first expression is the geometric cross section for
interacting at all, which is \signd in these units. The enhancement factor
$f(b)$ is the number that multiplies every rate in the rest of the chapter.
Its definition is the one place where I first guessed wrong: it is not
$O(b)$ divided by its average over interacting collisions, which would give
$\langle f\rangle=1$ for minimum bias by a different route, but the ratio
above, and the two differ by a few percent at the $b$ of our event. Which
$b$ a collision has is then drawn with a weight that depends on what is
asked: for minimum bias $\propto P_{\mathrm{int}}(b)\,\mathrm{d}^2b$, for a
collision known to contain a hard scattering
$\propto O(b)\,\exp[-f(b)\,\sigma(>Q_{\mathrm{hard}})/\signd]\,\mathrm{d}^2b$,
the exponential being the probability that no \emph{other} scattering above
the hard scale happened, which for our event is $0.99996$ and does not
matter.

\paragraph{The numbers.} The log prints the overlap table with $k=93.0$
so that $\langle n\rangle=8.974$, $\bmean_{\mathrm{ND}}=1.24$ in the
profile's own units, and $f(0)=4.52$:
\begin{center}\footnotesize
\begin{tabular}{rrrrr}\toprule
$b/\bmean$ & $O(b)$ & $f(b)$ & $P_{\mathrm{int}}(b)$ & $n(b)$ \\ \midrule
$0$ & $0.448$ & $4.52$ & $1.000$ & $41.7$ \\
$0.25$ & $0.390$ & $3.93$ & $1.000$ & $36.2$ \\
$0.50$ & $0.257$ & $2.59$ & $1.000$ & $23.9$ \\
$0.75$ & $0.130$ & $1.32$ & $1.000$ & $12.1$ \\
$1.00$ & $0.052$ & $0.52$ & $0.992$ & $4.8$ \\
$1.25$ & $0.017$ & $0.17$ & $0.784$ & $1.5$ \\
$1.50$ & $0.004$ & $0.04$ & $0.331$ & $0.4$ \\
$2.00$ & $0.0002$ & $0.002$ & $0.016$ & $0.02$ \\
\bottomrule
\end{tabular}
\end{center}
Averaged with the minimum-bias weight $f$ gives $0.97$, with the
hard-process weight $2.22$; \PYTHIA measures $2.26$ over its two thousand
events. Our event is at $b=0.3136\bmean$, where the script finds $f=3.627$
against \PYTHIA's $3.709$ for the same $b$ in the \CMSSW log, a $2\%$
difference from the coarser integration of the script; the expected number
of secondary scatterings is $f\langle n\rangle=32.5$, and the \CMSSW event
got $31$, the standalone dressing $24$. Seed~1 with $b$ drawn lands at
$1.21\bmean$, $f=0.20$, $n(b)=1.8$ and one secondary scattering; nature's
events are spread over the whole of panel~(c) of
\cref{fig:uninvited_cousins:overlap_plot}, and that spread, not the
scattering process, is the first reason why the underlying event of two
events with identical hard processes can differ by a factor of ten.

\begin{figure}[htbp]\centering
\includegraphics[page=2,width=\textwidth]{MPI/Standalone/theory_multiparton_interactions.pdf}
\caption[The proton profile, the overlap and the impact-parameter
distributions]{Page~2 of \file{theory_multiparton_interactions.pdf}. (a) The
double-Gaussian matter density of CP5 against a single Gaussian. (b) The
overlap $O(b)$, the interaction probability $P_{\mathrm{int}}(b)$ and the
enhancement $f(b)$ of \cref{eq:uninvited_cousins:fb}. (c) The distribution
of $b/\bmean$ for minimum-bias collisions and for collisions with a hard
process, the latter against the $400$ dressings with $b$ drawn and
against \PYTHIA; our event and the \CMSSW event are the vertical line.}
\label{fig:uninvited_cousins:overlap_plot}
\end{figure}

\subsection{The chain of interactions below the hard scale}
\label{sec:uninvited_cousins:chain}

How are the $n(b)$ scatterings generated? Not all at once. If every
scattering took its partons from the full proton, the momentum sum could be
violated; the scatterings must be ordered and each must see what the
previous ones have taken. \PYTHIA orders them in \pt, hardest first, exactly
as the shower orders its emissions, and generates the next one with the
same veto algorithm (\cref{sec:mother_gives_birth:veto}). The hard process
sits at the top; the probability density for the next scattering at
$\pt$, given that the previous one was at $\pt^{\mathrm{prev}}$ and none
happened in between, is
\begin{equation}
  \frac{\mathrm{d}P}{\mathrm{d}\pt^2}=\frac{f(b)}{\signd}\,\frac{\mathrm{d}\sigma}{\mathrm{d}\pt^2}\;
  \exp\!\left[-\frac{f(b)}{\signd}\int_{\pt^2}^{(\pt^{\mathrm{prev}})^2}\frac{\mathrm{d}\sigma}{\mathrm{d}\pt'^2}\,\mathrm{d}\pt'^2\right],
  \label{eq:uninvited_cousins:chain}
\end{equation}
the Sudakov form factor of \cref{eq:mother_gives_birth:sudakov} with the
emission rate replaced by the scattering rate per non-diffractive collision.
The chain stops when the next scale falls below $\pt^{\min}$. The
overestimate that the veto algorithm needs is the regulator itself: with
$g(\pt^2)=C/(\pt^2+\ptzero^2)^2\ge\mathrm{d}\sigma/\mathrm{d}\pt^2$ the
integral is elementary and the trial scale is
\begin{equation}
  \pt^2=\left[\frac{1}{(\pt^{\mathrm{prev}})^2+\ptzero^2}-\frac{\signd\ln R}{f(b)\,C}\right]^{-1}-\ptzero^2,
  \qquad R\in(0,1),
  \label{eq:uninvited_cousins:trial}
\end{equation}
accepted with probability $(\mathrm{d}\sigma/\mathrm{d}\pt^2)/g(\pt^2)$.
Once a scale is accepted the rapidities $y_3$, $y_4$ are drawn flat within
the kinematic limits, which fixes $x_1$, $x_2$ by
\cref{eq:uninvited_cousins:dsigma}, the flavours and the channel are
chosen according to their weights at that point, the colour flow of
\PYTHIA's process tables is attached, and the system is showered from its
own \pt downwards with the shower of \cref{ch:mother_gives_birth}, with one
difference: the backward evolution of each side now reads the
\emph{depleted} density of the next subsection.

\begin{figure}[htbp]\centering
\resizebox{\textwidth}{!}{% ---------------------------------------------------------------------------
% The chain of interactions of our event: a ladder in pT from the hard
% scale 113.3 GeV down to pTmin = 0.2 GeV, with the 24 secondary interactions
% of the -v log ("the chain of interactions" section) as rungs, and the
% no-interaction probability between rungs written as a Sudakov-like factor.
% The vertical position is log10(pT). Right: the first rung worked out.
% ---------------------------------------------------------------------------
\begin{tikzpicture}[x=1cm,y=1cm,
  every node/.style={font=\footnotesize},
  rung/.style={line width=1.4pt,teal!80!black},
  hardrung/.style={line width=2.2pt,red!80!black},
  lab/.style={font=\scriptsize,align=left},
  head/.style={font=\small\bfseries},
  numbox/.style={draw,rounded corners=3pt,thin,fill=blue!6,align=left,inner sep=4pt,text width=8.8cm,font=\scriptsize}]
% y(pT) = 2.3*(log10(pT) + 0.7): pTmin 0.2 -> 0, 1 GeV -> 1.61, 10 -> 3.91, 113 -> 6.33
\node[head,anchor=west] at (-1.6,7.6) {(a) the ladder of our event};
\draw[-{Stealth}] (0,-0.3) -- (0,6.9) node[right,font=\scriptsize] {$\pt$ of the scattering [GeV]};
\foreach \p/\l in {0.2/0.2,0.5/0.5,1/1,2/2,5/5,10/10,20/20,50/50,100/100}
  {\draw (-0.08,{2.3*(log10(\p)+0.7)}) -- (0.08,{2.3*(log10(\p)+0.7)}) node[left=3pt,font=\scriptsize] {\l};}
\draw[hardrung] (0.3,{2.3*(log10(113.32)+0.7)}) -- (2.6,{2.3*(log10(113.32)+0.7)}) node[right,font=\scriptsize,text=red!80!black] {the hard process, $113.3\GeV$ (sets $\pt^{\max}$)};
\foreach \p in {11.453,7.679,4.474,3.651,3.418,3.141,2.800,2.191,2.183,1.788,1.757,1.629,1.499,1.430,1.315,1.142,1.095,0.896,0.816,0.706,0.603,0.551,0.487,0.434}
  {\draw[rung] (0.3,{2.3*(log10(\p)+0.7)}) -- (2.6,{2.3*(log10(\p)+0.7)});}
\node[lab,anchor=west,text=teal!80!black] at (2.8,{2.3*(log10(11.453)+0.7)}) {1: $11.45\GeV$, $gg\to gg$};
\node[lab,anchor=west,text=teal!80!black] at (2.8,{2.3*(log10(7.679)+0.7)}) {2: $7.68\GeV$, $gg\to gg$};
\node[lab,anchor=west,text=teal!80!black] at (2.8,{2.3*(log10(4.474)+0.7)}) {3: $4.47\GeV$, $ug\to ug$};
\node[lab,anchor=west,text=teal!80!black] at (2.8,{2.3*(log10(2.191)+0.7)+0.12}) {8, 9: $2.19$, $2.18\GeV$};
\node[lab,anchor=west,text=teal!80!black] at (2.8,{2.3*(log10(1.095)+0.7)-0.05}) {17 rungs between $1.8$ and $0.43\GeV$};
\node[lab,anchor=west,text=teal!80!black] at (2.8,{2.3*(log10(0.434)+0.7)-0.25}) {24: $0.43\GeV$; then nothing above $\pt^{\min}$};
\draw[dotted,black!60] (-0.1,0) -- (2.6,0) node[right,font=\scriptsize,text=black!65] {$\pt^{\min}=0.2\GeV$: the chain stops};
% the gap between the hard scale and rung 1
\draw[{Stealth}-{Stealth},gray!70] (1.45,{2.3*(log10(113.32)+0.7)-0.1}) -- (1.45,{2.3*(log10(11.453)+0.7)+0.1});
\node[lab,anchor=west,text=black!65,text width=4.3cm] at (0.35,{0.5*(2.3*(log10(113.32)+0.7)+2.3*(log10(11.453)+0.7))+0.15}) {no scattering between $113$ and $11.5\GeV$:\\$\exp[-f\,\sigma(>11.45)/\sigma_{\mathrm{ND}}]$};
% ---------------- (b) the numbers of rung 1 -------------------------------------------------
\node[head,anchor=west] at (8.6,7.6) {(b) the first rung, from the \code{-v} log};
\node[numbox,anchor=north west] at (8.6,6.9) {\textbf{evolve down from $113.32\GeV$}: trial scales from the overestimate $C/(\pt^2+p_{\mathrm{T}0}^2)^2$, accepted at $\pt=11.453\GeV$ with veto weight $0.174$};
\node[numbox,anchor=north west] at (8.6,5.8) {\textbf{kinematics}: $y_3=-2.787$, $y_4=-0.658$, so $x_1=(\pt/E_{\mathrm{cm}})(e^{y_3}+e^{y_4})=4.88\times10^{-4}$, $x_2=1.53\times10^{-2}$; $\sqrt{\shat}=37.1\GeV$, $\hat t=-1233$, $\hat u=-147\GeV^2$};
\node[numbox,anchor=north west] at (8.6,4.5) {\textbf{coupling and regulator}: $\alphas(\pt^2+p_{\mathrm{T}0}^2=133.3)=0.1725$; $\bigl(\pt^2/(\pt^2+p_{\mathrm{T}0}^2)\bigr)^2=0.969$};
\node[numbox,anchor=north west] at (8.6,3.45) {\textbf{depleted protons}: A has $X=0.983$ left, $x/X=4.96\times10^{-4}$, $xg=21.4$; B has $X=0.530$ left (its valence $u$ is gone), $x/X=0.0289$, $xg=4.26$};
\node[numbox,anchor=north west] at (8.6,2.2) {\textbf{channel}: $gg\to gg$ $65.1\%$, $gd\to dg$ $3.6\%$, $gu\to ug$ $3.2\%$, \ldots\ $\to$ $gg\to gg$ chosen; colour flow in A $(136,137)$, in B $(137,138)$, out $(139,138)$, $(136,139)$};
\node[numbox,anchor=north west] at (8.6,0.95) {\textbf{then}: showered from its own $\pt=11.45\GeV$: 3 ISR + 10 FSR, 15 final partons; the protons are depleted again and the next rung is drawn below $11.45\GeV$};
\end{tikzpicture}
}
\caption[The ladder of secondary scatterings and the first rung worked
out]{Left: the \pt ladder of our event, the hard process at $113.3\GeV$ and
the 24 secondary scatterings of the \code{-v} log as rungs on a logarithmic
scale; the gap between the hard scale and the first rung is the exponential
of \cref{eq:uninvited_cousins:chain}. Right: every number of the first rung,
from the trial scale and its veto weight to the showered system.}
\label{fig:uninvited_cousins:chain_sketch}
\end{figure}

\paragraph{The numbers.} With $f(b)=3.627$ the rate of
\cref{eq:uninvited_cousins:chain} integrates to $32.55$ expected secondary
scatterings below $113.3\GeV$. The first one is accepted at $11.453\GeV$
with veto weight $0.174$, after a quiet stretch whose probability,
$\exp[-3.627\times3.6/55.42]=0.79$, is not small: the rate above $10\GeV$ is
low. Its rapidities $y_3=-2.787$ and $y_4=-0.658$ give
$x_1=4.88\times10^{-4}$ and $x_2=1.53\times10^{-2}$, $\sqrt{\shat}=37.1\GeV$,
and the channel weights are $gg\to gg$ $65.1\%$, $gd\to dg$ $3.6\%$, $gu\to
ug$ $3.2\%$ and so on; $gg\to gg$ is chosen, with colour tags $(136,137)$
and $(137,138)$ in and $(139,138)$, $(136,139)$ out, and the system is
showered with 3 initial-state and 10 final-state branchings into 15 partons.
The second rung is unusual and worth a look: accepted at $7.679\GeV$ with
veto weight $0.020$, it has $y_4=+6.21$, so $x_1=0.281$, a hard gluon from
proton~A, and $x_2=8.8\times10^{-4}$, and its outgoing parton at $y=6.2$
carries $1909\GeV$ almost along the beam. Such a scattering is invisible to
any detector and still takes $28\%$ of proton~A. The ladder then descends
through $4.47$, $3.65$, $3.42$, $3.14$, $2.80$, $2.19$ and $2.18\GeV$ and
seventeen rungs between $1.8$ and $0.43\GeV$; after the 24th the next
trial falls below $\pt^{\min}$ and the chain stops. The log counted $303$
trial scales for $24$ acceptances: the overestimate is loose by an order
of magnitude at low \pt, where the depleted densities and the halved
channels pull the true rate down. The \CMSSW event's ladder, read from the
status-33 pairs of its record, runs from $6.65$ to $0.24\GeV$ with 31
rungs; the two ladders are compared in panel~(b) of
\cref{fig:uninvited_cousins:chain_plot}, and the multiplicity distributions
of the ensembles in panel~(c).

\begin{figure}[htbp]\centering
\includegraphics[page=3,width=\textwidth]{MPI/Standalone/theory_multiparton_interactions.pdf}
\caption[The no-interaction probability, the ladder and the
multiplicity]{Page~3 of \file{theory_multiparton_interactions.pdf}. (a) The
probability of no scattering between the hard scale and \pt,
\cref{eq:uninvited_cousins:chain}, for $f=1$ and for our $f=3.63$, with the
24 rungs of our event as vertical lines. (b) The ladder of our event against
the 31 secondary scatterings of the \CMSSW event. (c) The number of
scatterings per event, including the hard one, for the two standalone
ensembles and the \PYTHIA ensemble at fixed $b$, against a Poisson of the
expected mean: the standalone chain is $20\%$ short of \PYTHIA's, for
the reason given in \cref{sec:uninvited_cousins:caveats}.}
\label{fig:uninvited_cousins:chain_plot}
\end{figure}

\subsection{Which partons: flavours, momentum fractions and the depleted proton}
\label{sec:uninvited_cousins:flavours}

Guess first: after the hard process took a valence $u$ quark with $x=0.47$
from proton~B, what is the probability that the next scattering takes
another valence $u$ from the same proton? With the full proton it would be
whatever the PDF says; with half the proton gone and one of its two $u$
quarks gone, it must be lower, and after both $u$ quarks are gone it must
be zero. The depleted proton of Sj\"ostrand and Skands~\cite{Sjostrand:2004pf}
is the bookkeeping that makes this true. Three rules modify the density
seen by each new scattering, and the \code{Beam} class of the script is a
port of \PYTHIA's \code{BeamParticle} that applies them.

\begin{figure}[htbp]\centering
\resizebox{\textwidth}{!}{% ---------------------------------------------------------------------------
% The protons are used up. Three snapshots of proton A of our event (from the
% -v log): before the collision, after the hard process took a gluon
% (x = 0.0171), and after the 24 secondary scatterings (X = 0.110 left, one
% valence u gone, companions of the sea quarks taken). Bottom: the three
% things that happen to a density when a parton leaves: rescaling, valence
% counting, and a companion for every sea quark.
% ---------------------------------------------------------------------------
\begin{tikzpicture}[x=1cm,y=1cm,
  every node/.style={font=\footnotesize},
  proton/.style={draw,circle,very thick,fill=gray!8,minimum size=30mm},
  val/.style={circle,fill=blue!70!black,inner sep=0pt,minimum size=3.2mm,text=white,font=\tiny},
  gone/.style={circle,draw=blue!70!black,dashed,thick,fill=white,inner sep=0pt,minimum size=3.2mm,text=blue!70!black,font=\tiny},
  sea/.style={circle,fill=violet!70,inner sep=0pt,minimum size=1.8mm},
  seabar/.style={circle,draw=violet!70,thick,fill=white,inner sep=0pt,minimum size=1.8mm},
  gluon/.style={thick,decorate,decoration={coil,aspect=0,segment length=3pt,amplitude=1.2pt},gray!70},
  bar/.style={draw,thin},
  lab/.style={font=\scriptsize,align=center},
  head/.style={font=\small\bfseries},
  bigarrow/.style={-{Stealth[length=3mm,width=3mm]},gray!60,line width=2.4pt}]
% ---------------- three snapshots of proton A --------------------------------------------------
\node[head,anchor=west] at (-2.2,2.6) {(a) proton A of our event, used up scattering by scattering};
\foreach \cx/\X/\title in {0/1.000/{before: $X=1$}, 5.0/0.983/{after the hard process: $X=0.983$}, 10.0/0.110/{after 24 more: $X=0.110$}} {
  \node[proton] at (\cx,0) {};
  \node[lab,anchor=south] at (\cx,1.65) {\title};
  % momentum bar below
  \draw[bar] (\cx-1.4,-2.05) rectangle (\cx+1.4,-1.75);
  \fill[teal!60] (\cx-1.4,-2.05) rectangle (\cx-1.4+2.8*\X,-1.75);
  \node[font=\scriptsize,anchor=north] at (\cx,-2.1) {momentum fraction left $X=1-\sum x_i$};}
% snapshot 1: 2 u + 1 d valence, gluons, sea pairs
\node[val] at (-0.5,0.4) {$u$}; \node[val] at (0.55,0.25) {$u$}; \node[val] at (0.05,-0.6) {$d$};
\draw[gluon] (-0.5,0.4) -- (0.55,0.25); \draw[gluon] (0.55,0.25) -- (0.05,-0.6); \draw[gluon] (0.05,-0.6) -- (-0.5,0.4);
\draw[gluon] (-0.5,0.4) -- (-1.1,0.9); \draw[gluon] (0.55,0.25) -- (1.15,-0.15);
\node[sea] at (-0.9,-0.4) {}; \node[seabar] at (-0.7,-0.75) {};
\node[sea] at (0.8,0.95) {}; \node[seabar] at (1.0,0.7) {};
% snapshot 2: a gluon gone (dashed), the rest rescaled
\node[val] at (4.5,0.4) {$u$}; \node[val] at (5.55,0.25) {$u$}; \node[val] at (5.05,-0.6) {$d$};
\draw[gluon] (4.5,0.4) -- (5.55,0.25); \draw[gluon] (5.55,0.25) -- (5.05,-0.6); \draw[gluon] (5.05,-0.6) -- (4.5,0.4);
\draw[gluon,red!70!black,dashed] (5.55,0.25) -- (6.15,-0.15);
\node[font=\scriptsize,text=red!70!black,anchor=west] at (6.2,-0.2) {$g$, $x=0.017$};
\node[sea] at (4.1,-0.4) {}; \node[seabar] at (4.3,-0.75) {};
\node[sea] at (5.8,0.95) {}; \node[seabar] at (6.0,0.7) {};
% snapshot 3: one valence u gone, sea quarks gone and their companions left, gluons gone
\node[gone] at (9.5,0.4) {$u$}; \node[val] at (10.55,0.25) {$u$}; \node[val] at (10.05,-0.6) {$d$};
\draw[gluon] (10.55,0.25) -- (10.05,-0.6);
\node[seabar] at (9.1,-0.4) {}; \node[seabar] at (9.3,-0.75) {}; \node[sea] at (10.8,0.95) {}; \node[sea] at (9.6,-1.0) {}; \node[seabar] at (10.9,-0.5) {};
\node[lab,anchor=west,text=black!65,text width=3.2cm] at (11.6,0.3) {25 partons given away: 15 gluons, 1 valence $u$, 9 sea quarks; left: $ud_0$ diquark and 6 companions};
% ---------------- (b) the three rules --------------------------------------------------------
\begin{scope}[yshift=-3.6cm]
\node[head,anchor=west] at (-2.2,0.0) {(b) what one departure does to the density of what remains};
\node[lab,anchor=north west,text width=4.0cm,align=left] at (-2.2,-0.4) {\textbf{rescale.} The others share the remaining $X$: a parton at $x$ is drawn from $f(x/X)/X$ (our rung 1: $x/X=4.96\times10^{-4}$ for A).};
\node[lab,anchor=north west,text width=4.0cm,align=left] at (2.6,-0.4) {\textbf{count valence.} Each valence quark can leave once: the $u$ that B gave the hard process is gone for good, so B keeps one $u$ and the $d$, and its valence-$u$ density is halved.};
\node[lab,anchor=north west,text width=4.3cm,align=left] at (7.4,-0.4) {\textbf{leave a companion.} A sea quark came from $g\to q\bar q$, so its partner stays behind, with $x_c$ drawn from the $g\to q\bar q$ kernel weighted by the gluon shape $(1-x_g)^4/x_g$, $x_g=x_c+x_s$.};
\begin{scope}[xshift=12.3cm,yshift=-0.9cm]
\draw[gluon] (0,0) -- (0.9,0);
\draw[thick,violet!80!black,-{Stealth[length=1.6mm]}] (0.9,0) -- (1.9,0.45) node[right,font=\scriptsize] {$q$ leaves, $x_s$};
\draw[thick,violet!80!black,-{Stealth[length=1.6mm]}] (0.9,0) -- (1.9,-0.45) node[right,font=\scriptsize] {$\bar q$ stays, $x_c$};
\end{scope}
\end{scope}
\end{tikzpicture}
}
\caption[The protons are used up]{Top: proton A of our event before the
collision, after the hard process took a gluon with $x=0.017$, and after the
24 secondary scatterings, with the momentum fraction $X$ left in each case;
a dashed circle is a valence quark that has left, hollow violet dots are the
companions owed to sea quarks that have left. Bottom: the three rules of
\cref{eq:uninvited_cousins:depleted} that modify the density seen by the
next scattering.}
\label{fig:uninvited_cousins:depletion_sketch}
\end{figure}

Let $X=1-\sum_i x_i$ be the momentum fraction not yet given away by a
proton. The modified density for a parton of flavour $q$ at fraction $x$ is
\begin{equation}
  xf_q^{\mathrm{mod}}(x)=\frac{1}{X}\left[\,
    \frac{n_q^{\mathrm{left}}}{n_q^{\mathrm{val}}}\;x'f_q^{\mathrm{val}}(x')
    +r_{\mathrm{gs}}\;x'f_q^{\mathrm{sea}}(x')
    +\sum_{\text{companions }c}x'f_c(x';x_s)\right],
  \qquad x'=\frac{x}{X},
  \label{eq:uninvited_cousins:depleted}
\end{equation}
with three ingredients. \emph{Rescaling:} every density is evaluated at
$x/X$ and divided by $X$, so that the remaining partons share the remaining
momentum. \emph{Valence counting:} a proton has two $u$ and one $d$ valence
quarks; each can leave once, and the valence density is scaled by the
number left. \emph{Companions:} a sea quark that leaves was one half of a
$g\to q\bar q$ splitting, so its partner is still in the proton with a
density derived from the splitting kernel,
$f_c(x_c;x_s)\propto g(x_c+x_s)\,P_{g\to q\bar q}\bigl(x_s/(x_c+x_s)\bigr)/(x_c+x_s)$
with $g(x)\propto(1-x)^4/x$, and it must leave with the remnant if nothing
takes it first. The gluon and sea densities are rescaled together by
$r_{\mathrm{gs}}=(1-\langle x\rangle_{\mathrm{val,left}})/(1-\langle x\rangle_{\mathrm{val,all}})$,
the momentum the departed valence quarks were carrying on average, so
that the momentum sum rule holds for the proton that remains.

\paragraph{The numbers.} After the hard process the log says that proton A
gave a gluon at $x=0.01714$ and has $X=0.9829$ left with both $u$ and the
$d$ still inside, and proton B gave the valence $u$ at $x=0.470$ and has
$X=0.530$, one $u$ and one $d$ left, with its gluon and sea rescaled by
$r_{\mathrm{gs}}=1.164$. At the second rung, where proton A is asked for
$x_1=0.281$, that is $x/X=0.287$, the modified densities are $xg=0.179$,
$xu=0.492$, $xd=0.209$, and for the first time a quark channel leads the
weights, $ug\to gu$ with $35.4\%$ against $gg\to gg$ with $28.9\%$; the
random number chose the gluons. At the third rung proton A supplies a $u$
quark that the log labels ``sea (companion still inside)'': from that
moment proton A owes a $\bar u$ companion, and at the end of the chain it
has given away $25$ partons, $15$ gluons, one valence $u$ and $9$ sea
quarks, with $X=0.110$ left, one $u$ and the $d$ for the diquark, and six
companions to pay. Proton B gave $25$ partons as well and keeps
$X=0.266$. Panel~(c) of \cref{fig:uninvited_cousins:flavours_plot} is the
staircase of $X$ for both protons: the second rung is the big step of
proton~A, the hard process that of proton~B.

\begin{figure}[htbp]\centering
\includegraphics[page=4,width=\textwidth]{MPI/Standalone/theory_multiparton_interactions.pdf}
\caption[Momentum fractions, processes and the depletion of the
protons]{Page~4 of \file{theory_multiparton_interactions.pdf}. (a) The
gluon, valence and sea densities at $Q=1.65\GeV$ with the $(x_1,x_2)$ of the
hard process and of the 24 secondary scatterings of our event. (b) The
processes of the secondary scatterings: $gg\to gg$ in all but five. (c) The
momentum fraction $X$ left in each proton after each scattering,
\cref{eq:uninvited_cousins:depleted}.}
\label{fig:uninvited_cousins:flavours_plot}
\end{figure}

\subsection{What is left: beam remnants and primordial transverse momentum}
\label{sec:uninvited_cousins:remnants}

The record of \cref{ch:mother_gives_birth} ended with a momentum sum that
was not the protons'. Here it becomes so. What is left of each proton after
the chain, $X=0.110$ of proton A and $0.266$ of proton B, must be given to
partons that hadronisation can handle, and the systems themselves must be
given the small transverse momentum that the partons had inside the
proton. Both are the business of \PYTHIA's \code{BeamRemnants}~\cite{Sjostrand:2004pf},
and the script ports its three steps.

\begin{figure}[htbp]\centering
\resizebox{\textwidth}{!}{% ---------------------------------------------------------------------------
% Beam remnants and primordial kT. Left: the light-cone momentum p+ of
% proton A shared between the hard system, the 24 secondary systems and the
% remnant partons (the -v log, "beam remnants" section: X = 0.110 left for A,
% 0.266 for B; remnants of A: ud_0 and six companions). Middle: the
% primordial-kT width against the scale of the system, with the hard system
% and the hardest MPI marked. Right: the kicks of one system and the recoil.
% ---------------------------------------------------------------------------
\begin{tikzpicture}[x=1cm,y=1cm,
  every node/.style={font=\footnotesize},
  lab/.style={font=\scriptsize,align=center},
  head/.style={font=\small\bfseries},
  kick/.style={-{Stealth[length=1.8mm]},thick}]
% ---------------- (a) light-cone bookkeeping ----------------------------------------------
\node[head,anchor=west] at (-0.4,3.3) {(a) the momentum must close};
\node[font=\scriptsize,anchor=west] at (-0.4,2.85) {$p^+_A=\sum_{\mathrm{systems}}p^+_i+\sum_{\mathrm{remnants}}p^+_r$, the same for $p^-_B$};
\draw[thin] (0,0) rectangle (4.4,1.0);
% hard system share of A: x = 0.0171 (tiny), MPI systems: 0.873, remnants 0.110
\fill[red!80!black] (0,0) rectangle (0.075,1.0);
\fill[teal!60] (0.075,0) rectangle (3.916,1.0);
\foreach \x in {0.3,0.55,1.9,2.25,2.9,3.2,3.45,3.6,3.72,3.8,3.86} {\draw[white,thin] (\x,0) -- (\x,1.0);}
\fill[gray!55] (3.916,0) rectangle (4.4,1.0);
\node[font=\scriptsize,text=red!80!black,anchor=south] at (0.1,1.05) {hard $x=0.017$};
\node[font=\scriptsize,text=teal!80!black] at (2.0,0.5) {24 secondary systems: $\sum x=0.873$};
\node[font=\scriptsize,text=white] at (4.16,0.5) {$0.11$};
\node[lab,anchor=north] at (2.2,-0.1) {proton A: $X=0.110$ is left for the remnants};
\draw[thin] (0,-2.4) rectangle (4.4,-1.4);
\fill[red!80!black] (0,-2.4) rectangle (2.068,-1.4);
\fill[teal!60] (2.068,-2.4) rectangle (3.231,-1.4);
\foreach \x in {2.3,2.5,2.75,2.95,3.1} {\draw[white,thin] (\x,-2.4) -- (\x,-1.4);}
\fill[gray!55] (3.231,-2.4) rectangle (4.4,-1.4);
\node[font=\scriptsize,text=white] at (1.03,-1.9) {hard: valence $u$, $x=0.47$};
\node[font=\scriptsize,text=white] at (3.8,-1.9) {$0.27$};
\node[lab,anchor=north] at (2.2,-2.5) {proton B: $X=0.266$ is left};
\node[lab,anchor=north west,text=black!65,align=left,text width=4.6cm] at (-0.3,-3.2) {remnants of A: $ud_0$ ($x=0.18$), $u$, $u$, $\bar u$, $\bar s$, $s$, $\bar u$ (companions);\\ of B: $ud_1$, $u$, $u$, $\bar s$, $u$, $\bar s$, $\bar u$, $\bar s$, $d$. Each system keeps its mass and rapidity; the remnants take $p^+=1420$, $p^-=3593\GeV$.};
% ---------------- (b) the primordial-kT width ---------------------------------------------
\begin{scope}[xshift=6.2cm,yshift=-2.0cm]
\node[head,anchor=west] at (-0.3,5.3) {(b) how hard the kick};
\draw[-{Stealth}] (0,0) -- (5.2,0) node[right] {$Q$ [GeV]};
\draw[-{Stealth}] (0,0) -- (0,4.6) node[above,font=\scriptsize] {$\sigma_{k_{\mathrm{T}}}$ [GeV]};
% x = 1.5*(log10(Q)+1): Q=0.1 -> 0, 1 -> 1.5, 10 -> 3, 100 -> 4.5 ; y = 2.2*sigma
\foreach \q/\l in {0.1/0.1,1/1,10/10,100/100} {\draw ({1.5*(log10(\q)+1)},0) -- ({1.5*(log10(\q)+1)},-0.1) node[below] {\l};}
\foreach \s/\l in {0.9/0.9,1.8/1.8} {\draw (0,{2.2*\s}) -- (-0.1,{2.2*\s}) node[left] {\l};}
\draw[very thick,teal!80!black] plot[domain=-1:2.1,samples=80,variable=\L] ({1.5*(\L+1)},{2.2*((1.5*0.9+pow(10,\L)*1.8)/(1.5+pow(10,\L)))});
\draw[dotted,black!60] (0,{2.2*0.9}) -- (5,{2.2*0.9}); \draw[dotted,black!60] (0,{2.2*1.8}) -- (5,{2.2*1.8});
\node[font=\scriptsize,text=teal!80!black,anchor=north west,text width=4.6cm] at (0.15,1.75) {$\sigma=\dfrac{Q_{1/2}\,\sigma_{\mathrm{soft}}+Q\,\sigma_{\mathrm{hard}}}{Q_{1/2}+Q}$, $Q_{1/2}=1.5$, $\sigma_{\mathrm{soft}}=0.9$, $\sigma_{\mathrm{hard}}=1.8\GeV$; remnants: $0.4\GeV$ flat};
\node[circle,fill=red!80!black,inner sep=1.8pt] at ({1.5*(log10(113.3)+1)},{2.2*1.786}) {};
\node[font=\scriptsize,text=red!80!black,anchor=north] at ({1.5*(log10(113.3)+1)},{2.2*1.786-0.15}) {hard: $Q=113$, $1.79\GeV$};
\node[circle,fill=teal!80!black,inner sep=1.8pt] at ({1.5*(log10(11.45)+1)},{2.2*1.676}) {};
\node[font=\scriptsize,text=teal!80!black,anchor=north east] at ({1.5*(log10(11.45)+1)-0.1},{2.2*1.676-0.1}) {rung 1: $1.68$};
\node[circle,fill=teal!80!black,inner sep=1.8pt] at ({1.5*(log10(0.434)+1)},{2.2*1.10}) {};
\node[font=\scriptsize,text=teal!80!black,anchor=north west] at ({1.5*(log10(0.434)+1)+0.05},{2.2*1.10-0.12}) {rung 24: $1.10$};
\end{scope}
% ---------------- (c) the kicks of the hard system --------------------------------------
\begin{scope}[xshift=14.2cm,yshift=0.2cm]
\node[head,anchor=west] at (-1.9,3.1) {(c) the kicks sum to zero per proton};
\draw[gray!50,thin] (-1.8,0) -- (1.8,0); \draw[gray!50,thin] (0,-1.8) -- (0,1.8);
\node[font=\scriptsize,text=black!65] at (1.75,-0.2) {$k_x$}; \node[font=\scriptsize,text=black!65] at (0.2,1.75) {$k_y$};
% hard system initiators: A (-1.767,+1.652), B (-1.113,-1.838); scale 0.5 cm/GeV
\draw[kick,red!80!black] (0,0) -- (-0.884,0.826) node[above,font=\scriptsize] {A, hard: $(-1.77,+1.65)$};
\draw[kick,red!80!black,dashed] (0,0) -- (-0.557,-0.919) node[below,font=\scriptsize] {B, hard: $(-1.11,-1.84)$};
\draw[kick,teal!80!black] (0,0) -- (-0.449,0.506);
\draw[kick,teal!80!black] (0,0) -- (0.533,-0.555);
\draw[kick,teal!80!black] (0,0) -- (0.051,-0.327);
\node[font=\scriptsize,text=teal!80!black,anchor=west] at (0.6,-0.5) {rungs 1--3 (A side)};
\draw[kick,gray!70!black] (0,0) -- (0.55,0.3);
\node[font=\scriptsize,text=black!65,anchor=west] at (0.25,0.55) {remnants: what balances};
\node[lab,anchor=north,text=black!65] at (0,-2.9) {every initiator and remnant of a proton gets a Gaussian $\vec k_{\mathrm{T}}$;\\the imbalance is shared out so that $\sum\vec k_{\mathrm{T}}=0$ per proton};
\end{scope}
\end{tikzpicture}
}
\caption[Beam remnants and primordial transverse momentum]{Left: the
light-cone momentum of each proton shared between the hard system, the 24
secondary systems and the remnants, with the numbers of our event. Middle:
the primordial-$k_{\mathrm{T}}$ width of \cref{eq:uninvited_cousins:kt}
against the scale of the system, with the hard system and two rungs of our
event. Right: the kicks of the hard system's two initiators and the recoil
that the rest of the proton absorbs.}
\label{fig:uninvited_cousins:remnants_sketch}
\end{figure}

\emph{Flavours.} The valence quarks left are paired into a diquark, two
valence quarks into $ud_0$ or $ud_1$ (spin 0 or 1, in the ratio $3:1$) or
$uu_1$; if a valence quark is left over it goes alone; every companion owed
is added; and a proton with nothing left at all gets a gluon. \emph{Momenta.}
Each remnant parton draws an unrescaled fraction from a shape chosen by its
kind, $(1-x)^{3.5}/\sqrt x$ for a valence $u$, $(1-x)^2/\sqrt x$ for a
valence $d$, the sum of two such draws for a diquark, $(1-x)^4/x$ for a
gluon, the companion distribution of \cref{eq:uninvited_cousins:depleted}
for a companion, and all remnant fractions are then rescaled so that they
sum to $X$. \emph{Transverse momentum.} Every initiator and every remnant
of a proton receives a Gaussian $\vec k_{\mathrm{T}}$ of width
\begin{equation}
  \sigma_{k_{\mathrm{T}}}(Q,m)=\frac{Q_{1/2}\,\sigma_{\mathrm{soft}}+Q\,\sigma_{\mathrm{hard}}}{Q_{1/2}+Q}\cdot\frac{m}{m+m_{1/2}},
  \label{eq:uninvited_cousins:kt}
\end{equation}
with $Q_{1/2}=1.5\GeV$, $\sigma_{\mathrm{soft}}=0.9\GeV$,
$\sigma_{\mathrm{hard}}=1.8\GeV$ and $m_{1/2}=1\GeV$ (the CP5 values of
\code{BeamRemnants:halfScaleForKT}, \code{primordialKTsoft},
\code{primordialKThard} and \code{halfMassForKT}), where $Q$ is the \pt of
the scattering (the renormalisation scale for the hard process) and $m$ the
mass of the system; the remnants get $0.4\GeV$ flat. The vector sum over a
proton is then shared out so that it vanishes. With the fractions and the transverse momenta fixed, each system
is required to keep its invariant mass and its rapidity, which determines
the light-cone momenta $p^+$ and $p^-$ of its two initiators; the remnants
receive what is left of $p^+_A$ and $p^-_B$, and the record's momentum sum
closes on the two protons. The script solves the light-cone conditions with
\PYTHIA's exact iteration and prints the remnant systems' masses.

\paragraph{The numbers.} Proton A's remnant is a $ud_0$ diquark of mass
$0.579\GeV$ with six companions, $u$, $u$, $\bar u$, $\bar s$, $s$, $\bar u$;
proton B's a $ud_1$ with eight, $u$, $u$, $\bar s$, $u$, $\bar s$, $\bar u$,
$\bar s$, $d$; sixteen remnant partons in all, status~63. For the hard
system, $Q=113.3\GeV$ and $m=1221\GeV$, the width is $1.786\GeV$ and the
kicks given are $(-1.767,+1.652)\GeV$ to the initiator from A and
$(-1.113,-1.838)\GeV$ to the one from B; for the first rung, $Q=11.45\GeV$,
$1.676\GeV$; for the last, $Q=0.43\GeV$, $1.10\GeV$. The light-cone
solution was found in one attempt and leaves $p^+=1419.9\GeV$ for the
remnants of A and $p^-=3593.3\GeV$ for those of B, remnant systems of mass
$7.2$ and $61.8\GeV$; the remnant fractions were rescaled by $1.585$ and
$1.480$ to fit. The $ud_0$ carries $x=0.184$, $206.6\GeV$ at $\eta=6.6$; the
hardest remnant of B is a $\bar s$ companion with $x=0.307$, $817\GeV$ at
$\eta=-8.4$; and the softest, a $\bar u$ of $0.67\GeV$, sits at $\eta=0.35$
in the middle of the detector, the one remnant that will be seen. The
hard system's kick is the first thing of this stage that touches the
ledger: boosted to the frame of its kicked initiators, the $b$ quark moves
from $78.12$ to $78.29\GeV$ and its family from $100.45$ to $100.67\GeV$.

\begin{figure}[htbp]\centering
\includegraphics[page=5,width=\textwidth]{MPI/Standalone/theory_multiparton_interactions.pdf}
\caption[Remnant momentum shapes, primordial $k_{\mathrm{T}}$ and the
light-cone bookkeeping]{Page~5 of \file{theory_multiparton_interactions.pdf}.
(a) The unrescaled momentum shapes of the remnant partons, with the sixteen
remnants of our event as markers. (b) The primordial-$k_{\mathrm{T}}$ width
of \cref{eq:uninvited_cousins:kt} against the scale, the width used for each
of the 25 systems of our event after the mass damping, and the mean kick
actually given. (c) The share of the light-cone momentum of each proton
taken by the hard system, the secondary systems and the remnants: the bars
close at one.}
\label{fig:uninvited_cousins:remnants_plot}
\end{figure}

\subsection{From formula to event: the stage step by step}
\label{sec:uninvited_cousins:step_by_step}

\Cref{fig:uninvited_cousins:flow} collects the construction in the order the
script runs it, with the numbers of our event beside each step. Read it
once now and once more after the next two sections; the colour
reconnection, step~7, is the only step that has not yet been explained.

\begin{figure}[htbp]\centering
\resizebox{\textwidth}{!}{% ---------------------------------------------------------------------------
% The stage as a flow chart: the steps of standalone_multiparton_interactions.py
% on the left, the numbers of our event (the -v log, seed 1, b = 0.3136 <b>)
% on the right. Function names in brackets.
% ---------------------------------------------------------------------------
\begin{tikzpicture}[x=1cm,y=1cm,
  every node/.style={font=\footnotesize},
  box/.style={draw,rounded corners=3pt,thick,fill=gray!8,align=left,inner sep=5pt,minimum height=10mm,text width=5.6cm},
  stepbox/.style={box,fill=teal!10},
  numbox/.style={draw,rounded corners=3pt,thin,fill=blue!6,align=left,inner sep=4pt,text width=8.6cm,font=\scriptsize},
  stepno/.style={draw,circle,inner sep=1pt,font=\scriptsize\bfseries,fill=teal!35,minimum size=5mm},
  flow/.style={-{Stealth[length=2.5mm,width=2mm]},thick},
  fn/.style={font=\scriptsize\ttfamily,text=black!65},
  head/.style={font=\small\bfseries}]
\node[head,anchor=west] at (-0.6,1.4) {the recipe};
\node[head,anchor=west] at (6.7,1.4) {our event (seed 1, $b=0.3136\,\langle b\rangle$)};
\node[stepbox] (S1) at (2.5,0) {\textbf{1\; set the scene}\\ $p_{\mathrm{T}0}(E_{\mathrm{cm}})$, the regularised $\mathrm{d}\sigma/\mathrm{d}\pt^2$ integrated above $\pt^{\min}$, $\sigma_{\mathrm{ND}}$, the overlap $O(b)$ and $k$};
\node[stepbox,below=7mm of S1] (S2) {\textbf{2\; draw the impact parameter}\\ with weight $O(b)$ for a hard event, then $f(b)=n(b)/\langle n\rangle$};
\node[stepbox,below=7mm of S2] (S3) {\textbf{3\; next interaction}\\ veto algorithm downwards in $\pt$ with rate $f(b)\,\mathrm{d}\sigma/\sigma_{\mathrm{ND}}$; stop below $\pt^{\min}$};
\node[stepbox,below=7mm of S3] (S4) {\textbf{4\; dress the scattering}\\ $y_3,y_4\to x_1,x_2$; flavours from the depleted densities; process, colour flow, kinematics; shower it (chapter 2)};
\node[stepbox,below=7mm of S4] (S5) {\textbf{5\; deplete both protons}\\ rescale $x$, count valence, add a companion for each sea quark; back to~3};
\node[stepbox,below=7mm of S5] (S6) {\textbf{6\; beam remnants}\\ flavours and diquark, $x$ shapes, primordial $k_{\mathrm{T}}$, light-cone solution for every system};
\node[stepbox,below=7mm of S6] (S7) {\textbf{7\; colour reconnection}\\ merge soft systems into harder ones, string the remnants, measure $\lambda$};
\node[stepbox,below=7mm of S7] (S8) {\textbf{8\; write the record}\\ status $-31$, $-33$, $-61$, $62$, $63$; check $\sum p$ and the colour lines};
\foreach \a/\b in {S1/S2,S2/S3,S3/S4,S4/S5,S5/S6,S6/S7,S7/S8} {\draw[flow] (\a) -- (\b);}
\draw[flow] (S5.west) -- ++(-0.9,0) |- (S3.west);
\foreach \n/\s in {1/S1,2/S2,3/S3,4/S4,5/S5,6/S6,7/S7,8/S8} {\node[stepno] at ($(\s.west)+(-0.35,0)$) {\n};}
\node[fn,anchor=south east] at (S1.north east) {SigmaInt, Overlap};
\node[fn,anchor=south east] at (S3.north east) {Generator.next\_pt};
\node[fn,anchor=south east] at (S4.north east) {Generator.scatter, Shower};
\node[fn,anchor=south east] at (S5.north east) {Beam};
\node[fn,anchor=south east] at (S6.north east) {remnants, set\_kinematics};
\node[fn,anchor=south east] at (S7.north east) {reconnect\_mpi, remnant\_colours};
\node[numbox,anchor=west] (N1) at (6.7,0) {$p_{\mathrm{T}0}=1.41\,(13600/7000)^{0.03344}=1.4417\GeV$; $\sigma_{\mathrm{int}}=497.3\,\mathrm{mb}$ (\CMSSW log: $500.67$);\\ $\sigma_{\mathrm{ND}}=100.309-21.89-23.0=55.42\,\mathrm{mb}$; $\langle n\rangle=8.97$; $k=93.0$, $f(0)=4.52$};
\node[numbox,anchor=west] (N2) at (6.7,0 |- S2) {fixed to the \CMSSW event: $b=0.3136\,\langle b\rangle$, $f(b)=3.627$ ($3.709$ in \CMSSW);\\ expected $32.5$ interactions below $\pt^{\max}=113.3\GeV$; seed 1 with $b$ drawn would give $b=1.21\,\langle b\rangle$, $f=0.20$};
\node[numbox,anchor=west] (N3) at (6.7,0 |- S3) {rung 1 at $11.453\GeV$ (weight $0.174$), 2 at $7.679$, 3 at $4.474$, \ldots, 24 at $0.434\GeV$;\\ 303 trial scales in all; after rung 24 nothing above $0.2\GeV$};
\node[numbox,anchor=west] (N4) at (6.7,0 |- S4) {rung 1: $y_3=-2.79$, $y_4=-0.66$, $x_1=4.9\times10^{-4}$, $x_2=0.015$, $gg\to gg$ ($65\%$ of the weight);\\ showered: 3 ISR + 10 FSR, 15 partons. All 24: 31 ISR + 76 FSR, 155 final partons};
\node[numbox,anchor=west] (N5) at (6.7,0 |- S5) {A: $X$ $0.983\to0.980\to0.672\to\cdots\to0.110$, one valence $u$ gone, 6 companions owed;\\ B: $0.530\to0.506\to\cdots\to0.266$ (the hard $u$ was valence), 8 companions owed};
\node[numbox,anchor=west] (N6) at (6.7,0 |- S6) {A: $ud_0+u,u,\bar u,\bar s,s,\bar u$; B: $ud_1+u,u,\bar s,u,\bar s,\bar u,\bar s,d$ (16 remnants);\\ $\sigma_{k_{\mathrm{T}}}=1.79\GeV$ for the hard system; remnant systems of mass $7.2$ and $61.8\GeV$; solved in 1 attempt};
\node[numbox,anchor=west] (N7) at (6.7,0 |- S7) {23 of 24 systems merged (rung 1 at $11.45\GeV$ survives, $P=0.30$); $\lambda$: $805.7\to403.5$; 51 strings;\\ 6 systems become colour-singlet gluon pairs};
\node[numbox,anchor=west] (N8) at (6.7,0 |- S8) {963 rows, 211 final partons (40 hard system + 155 secondary + 16 remnants);\\ $\sum p$ off $(13600,0,0,0)$ by $2\times10^{-3}\GeV$; written to \file{example/MPI/event_MPI.lhe}};
\foreach \s/\n in {S1/N1,S2/N2,S3/N3,S4/N4,S5/N5,S6/N6,S7/N7,S8/N8} {\draw[flow,gray!60] (\s.east) -- (\n.west);}
\end{tikzpicture}
}
\caption[The MPI stage as a flow chart]{The eight steps of
\file{standalone_multiparton_interactions.py}, with the names of the classes
and methods that implement them, and the numbers of our event at each step
from the \code{-v} log.}
\label{fig:uninvited_cousins:flow}
\end{figure}

\section{Colour reconnection and its effect on jet mass and top mass}
\label{sec:uninvited_cousins:colour_reconnection_and_its_effect_on_je}

Guess first: before hadronisation, the event has 25 systems, each with its
own colour lines running from its outgoing partons back through its
initiators to the remnants. If every system kept its own strings, how long
would the strings be? Very long: a $gg\to gg$ system at $y=0$ is connected
through its two initiators to remnants at $y=\pm8$, and every one of the
24 soft systems adds two such strings spanning the whole rapidity range.
Hadronisation turns string length into hadrons, so this naive picture
predicts a charged multiplicity that grows linearly with the number of
scatterings, and it does not: the multiplicity measured at the LHC grows
much more slowly~\cite{Sjostrand:2017cdm}. Something rearranges the colour
lines so that the total length is smaller. Nobody knows exactly what, and
this is one of the two subjects of the open problems at the end of the
chapter; \PYTHIA's default answer, the \emph{MPI-based} colour
reconnection~\cite{Sjostrand:2004pf,Sjostrand:2013cya}, is a simple rule
that the script ports.

\begin{figure}[htbp]\centering
\resizebox{\textwidth}{!}{% ---------------------------------------------------------------------------
% Colour reconnection in the MPI-based model. (a) Before: every system has
% its own colour strings, the hard system's chain and, separately, a soft
% gg -> gg system whose strings run far across rapidity. (b) After: the soft
% system's gluons are inserted into the dipole of the harder system that
% minimises (p_g.p_i)(p_g.p_j)/(p_i.p_j); the strings get shorter. (c) The
% reconnection probability with the systems of our event (the -v log,
% "colour reconnection" section): rung 1 at 11.45 GeV survives, rung 2 at
% 7.68 GeV is merged with P = 0.49, every softer rung with P > 0.74.
% ---------------------------------------------------------------------------
\begin{tikzpicture}[x=1cm,y=1cm,
  every node/.style={font=\footnotesize},
  q/.style={circle,inner sep=0pt,minimum size=2.4mm,fill=red!80!black},
  g/.style={circle,inner sep=0pt,minimum size=2.0mm,fill=green!50!black},
  gs/.style={circle,inner sep=0pt,minimum size=2.0mm,fill=teal!80!black},
  rem/.style={diamond,inner sep=0pt,minimum size=2.6mm,fill=gray!60},
  str/.style={line width=1.3pt,black!75},
  strs/.style={line width=1.3pt,teal!80!black},
  lab/.style={font=\scriptsize,align=center},
  head/.style={font=\small\bfseries}]
% ---------------- (a) before ----------------------------------------------------------------
\begin{scope}
\node[head,anchor=west] at (-0.5,3.2) {(a) before: one string per system};
\draw[-{Stealth}] (-0.3,-2.3) -- (5.4,-2.3) node[below left=2pt and -2pt] {$y$};
\draw[-{Stealth}] (-0.3,-2.3) -- (-0.3,2.8) node[above,font=\scriptsize] {$\phi$};
% hard system chain: remnant A ... g ... b ... g ... bbar ... remnant B (schematic)
\node[rem] (rA) at (5.2,0.6) {}; \node[g] (g1) at (3.9,1.9) {}; \node[q] (b) at (0.9,-1.5) {}; \node[g] (g2) at (2.0,0.2) {};
\node[g] (g3) at (-0.1,2.3) {}; \node[q] (bb) at (-0.1,1.0) {}; \node[rem] (rB) at (-0.2,-1.9) {};
\draw[str] (rA) -- (g1) -- (b) -- (g2) -- (bb) -- (g3) -- (rB);
\node[font=\scriptsize,text=red!80!black,anchor=north] at (b.south) {$b$};
\node[font=\scriptsize,text=red!80!black,anchor=south] at (bb.north) {$\bar b$};
% a soft gg -> gg system: its own two strings between the two outgoing gluons and the remnants
\node[gs] (s1) at (2.7,-1.6) {}; \node[gs] (s2) at (3.4,1.1) {};
\draw[strs] (rA) .. controls (4.4,-0.4) .. (s1);
\draw[strs] (s1) -- (s2);
\draw[strs] (s2) .. controls (1.4,2.6) .. (rB);
\node[lab,text=teal!80!black,anchor=west] at (3.55,-1.6) {soft $gg\to gg$\\$\pt=7.7\GeV$};
\node[lab,anchor=north,text=black!65] at (2.4,-2.6) {$\lambda=\sum_{\mathrm{pairs}}\ln(1+m_{ij}^2/m_0^2)=806$ before};
\end{scope}
% ---------------- (b) after -----------------------------------------------------------------
\begin{scope}[xshift=7.0cm]
\node[head,anchor=west] at (-0.5,3.2) {(b) after: soft gluons join harder dipoles};
\draw[-{Stealth}] (-0.3,-2.3) -- (5.4,-2.3) node[below left=2pt and -2pt] {$y$};
\draw[-{Stealth}] (-0.3,-2.3) -- (-0.3,2.8) node[above,font=\scriptsize] {$\phi$};
\node[rem] (rA) at (5.2,0.6) {}; \node[g] (g1) at (3.9,1.9) {}; \node[q] (b) at (0.9,-1.5) {}; \node[g] (g2) at (2.0,0.2) {};
\node[g] (g3) at (-0.1,2.3) {}; \node[q] (bb) at (-0.1,1.0) {}; \node[rem] (rB) at (-0.2,-1.9) {};
\node[gs] (s1) at (2.7,-1.6) {}; \node[gs] (s2) at (3.4,1.1) {};
\draw[str] (rA) -- (g1) -- (s2) -- (b) -- (g2) -- (bb) -- (g3) -- (rB);
\draw[strs,dashed] (g1) -- (b);
\node[font=\scriptsize,text=red!80!black,anchor=north] at (b.south) {$b$};
\node[font=\scriptsize,text=red!80!black,anchor=south] at (bb.north) {$\bar b$};
\draw[strs] (s1) .. controls (1.8,-1.0) .. (b);
\draw[strs] (s1) -- (g2);
\node[lab,text=teal!80!black,anchor=west,text width=2.2cm] at (3.55,-1.5) {each gluon goes where $\frac{(p_g\!\cdot p_i)(p_g\!\cdot p_j)}{p_i\cdot p_j}$ is smallest};
\node[lab,anchor=north,text=black!65] at (2.4,-2.6) {$\lambda=403$ after: $23$ of $24$ merged, $51$ strings};
\end{scope}
% ---------------- (c) the probability ------------------------------------------------------
\begin{scope}[xshift=14.2cm,yshift=-2.3cm]
\node[head,anchor=west] at (-0.3,5.5) {(c) who is merged};
\node[font=\scriptsize,anchor=west] at (-0.3,4.95) {$P=(R\,p_{\mathrm{T}0})^2/\bigl((R\,p_{\mathrm{T}0})^2+\pt^2\bigr)$};
\draw[-{Stealth}] (0,0) -- (5.0,0) node[right] {$\pt$ [GeV]};
\draw[-{Stealth}] (0,0) -- (0,4.4) node[above,font=\scriptsize] {$P$};
% x = 1.6*(log10(pT)+0.5): 0.316 -> 0, 1 -> 0.8, 10 -> 2.4, 100 -> 4.0 ; y = 4*P
\foreach \p/\l in {1/1,10/10,100/100} {\draw ({1.6*(log10(\p)+0.5)},0) -- ({1.6*(log10(\p)+0.5)},-0.1) node[below] {\l};}
\foreach \y/\l in {2/0.5,4/1} {\draw (0,\y) -- (-0.1,\y) node[left] {\l};}
\draw[very thick,teal!80!black] plot[domain=-0.5:2.4,samples=80,variable=\L] ({1.6*(\L+0.5)},{4/(1+pow(10,2*\L-1.7457))});
\draw[dashed,black!60] ({1.6*(log10(7.462)+0.5)},0) -- ({1.6*(log10(7.462)+0.5)},2.0);
\node[font=\scriptsize,text=black!65,anchor=west] at ({1.6*(log10(7.462)+0.5)+0.05},0.5) {$R\,p_{\mathrm{T}0}=5.176\times1.44=7.46\GeV$};
\node[circle,fill=red!80!black,inner sep=1.7pt] at ({1.6*(log10(11.453)+0.5)},{4*0.298}) {};
\node[font=\scriptsize,text=red!80!black,anchor=south west] at ({1.6*(log10(11.453)+0.5)},{4*0.298+0.05}) {rung 1: $P=0.30$, kept};
\node[circle,fill=teal!80!black,inner sep=1.7pt] at ({1.6*(log10(7.679)+0.5)},{4*0.4857}) {};
\node[font=\scriptsize,text=teal!80!black,anchor=south west] at ({1.6*(log10(7.679)+0.5)+0.05},{4*0.4857+0.1}) {rung 2: $0.49$, merged};
\node[circle,fill=teal!80!black,inner sep=1.7pt] at ({1.6*(log10(0.434)+0.5)},{4*0.9966}) {};
\node[font=\scriptsize,text=teal!80!black,anchor=north west] at ({1.6*(log10(0.434)+0.5)+0.1},{4*0.9966-0.1}) {rung 24: $0.997$};
\end{scope}
\end{tikzpicture}
}
\caption[Colour reconnection in the MPI-based model]{(a) Before
reconnection each system has its own strings; the hard system's chain runs
from one remnant through the $b$, its gluons and the $\bar b$ to the other,
and a soft $gg\to gg$ system adds two strings across the whole rapidity
range. (b) After reconnection the soft gluons have been inserted into the
dipoles of the harder system where
\cref{eq:uninvited_cousins:insertion} is smallest. (c) The merging
probability of \cref{eq:uninvited_cousins:cr_prob} with the rungs of our
event: only the first rung survives on its own.}
\label{fig:uninvited_cousins:cr_sketch}
\end{figure}

The rule has two parts. Each secondary system, starting from the softest,
is merged into a harder one with probability
\begin{equation}
  P_{\mathrm{merge}}=\frac{(R\,\ptzero)^2}{(R\,\ptzero)^2+\pt^2},
  \qquad R=\code{ColourReconnection:range}=5.176\ \text{(CP5)},
  \label{eq:uninvited_cousins:cr_prob}
\end{equation}
so that soft systems are almost always merged and systems well above
$R\,\ptzero$ almost never. When a system is merged, each of its final
gluons is inserted into the colour dipole $(i,j)$ of the harder system
that minimises
\begin{equation}
  \frac{(p_g\cdot p_i)(p_g\cdot p_j)}{p_i\cdot p_j},
  \label{eq:uninvited_cousins:insertion}
\end{equation}
the invariant that measures how much string length the insertion adds; a
$q\bar q$ pair from a gluon splitting in the merged system is inserted as
a pair. If both initiators of a merged system were gluons and all its
final partons have left, the two initiators form a colour-singlet pair and
the remnants see no string from that system at all. After the merging, the
remaining initiators and the remnants are strung together, one chain per
proton, with the resolution of the cases where both protons closed the
same colour line. The bookkeeping number that summarises the result is the
string length
\begin{equation}
  \lambda=\sum_{\text{adjacent pairs }(i,j)}\ln\!\left(1+\frac{m_{ij}^2}{m_0^2}\right),
  \qquad m_0=0.3\GeV,
  \label{eq:uninvited_cousins:lambda}
\end{equation}
which hadronisation will convert into hadrons at a rate of roughly one per
unit of $\lambda$.

\paragraph{The numbers.} $R\,\ptzero=5.176\times1.4417=7.462\GeV$. The log
tries the 24 systems from the softest upwards: rung 24 at $0.434\GeV$ has
$P=0.9966$ and is merged into rung 23, rung 23 into 22, and so on; rung 3
at $4.47\GeV$ has $P=0.736$ and is merged into rung 2; rung 2 at $7.68\GeV$
has $P=0.486$ and is merged, into the hard system; rung 1 at $11.45\GeV$
has $P=0.298$ and is not. $23$ of $24$ systems merged, and since the chain
of merges ends in the hard system, every soft gluon of the event ends up on
the hard system's dipoles. The log lists each insertion: the $7.34\GeV$
gluon of rung 2 goes into the dipole between rows 702 and 738, the two
initiators of the hard system, with the invariant $441.8\GeV^2$; the
$u\bar u$ pair of rung 2 into a dipole of the $\bar b$ family; the six
systems whose initiators were both gluons become colour-singlet
pass-throughs. The string length drops from $805.7$ to $403.5$ and the
event ends with $51$ colour-singlet strings. One of them passes through the
$b$ quark, drawn in panel~(c) of \cref{fig:uninvited_cousins:cr_plot}: it
now contains partons of secondary systems, cousins that hadronisation will
tie to the mother's family.

\paragraph{Why it matters for the jet mass and the top mass.} Colour
reconnection moves no momentum. Every number in
\cref{tab:uninvited_cousins:numbers} that is a momentum is the same with
\code{--no-cr}, and the ledger of this chapter does not see it. What it
changes is which partons will share a string, and therefore where the
hadrons go when the strings break in \cref{ch:children_getting_married}: a string
that connects the $b$ to a soft gluon at the other side of the event
produces hadrons between them, some inside the jet cone and some outside,
and the jet mass and the jet's energy are shifted by a few hundred MeV in
a way that no parton-level quantity shows. For a $b$ jet at $100\GeV$ that is
a fraction of a percent; for the top-quark mass, reconstructed from three
jets whose strings may or may not connect to each other and to the beam
remnants, Argyropoulos and Sj\"ostrand showed that different reconnection
models shift the measured mass by up to $0.5\GeV$~\cite{Argyropoulos:2014zoa},
which is the size of the total uncertainty of today's best
measurements~\cite{CMS:2023ebf,ATLAS:2024dxp}. CMS has tuned the
alternative models to underlying-event data~\cite{CMS:2022awf}, and the
spread between them is a systematic uncertainty in every top-mass result.

\begin{figure}[htbp]\centering
\includegraphics[page=6,width=\textwidth]{MPI/Standalone/theory_multiparton_interactions.pdf}
\caption[The reconnection probability, the string length and the string
through the $b$]{Page~6 of \file{theory_multiparton_interactions.pdf}.
(a) The merging probability of \cref{eq:uninvited_cousins:cr_prob} with
the 25 systems of our event, merged or not. (b) The string length
\cref{eq:uninvited_cousins:lambda} before and after reconnection. (c) The
string that passes through the $b$ quark after reconnection in the
$(y,\phi)$ plane, with its partons coloured by origin.}
\label{fig:uninvited_cousins:cr_plot}
\end{figure}

\section{Measuring the underlying event: transverse-region observables}
\label{sec:uninvited_cousins:measuring_the_underlying_event_transvers}

Guess first: the secondary scatterings of our event carry $230\GeV$ of
transverse momentum. How much of it would a measurement call ``underlying
event''? None of it directly, because no detector knows which parton came
from which scattering. What can be measured is the activity in a region of
the event where the hard process and its shower put little, and the
definition that CDF introduced~\cite{CDF:2001onq,Field:2002vt} and that
every LHC measurement uses~\cite{ATLAS:2017blj,CMS:2017ngy} divides the
azimuth relative to the leading object into three regions
(\cref{fig:uninvited_cousins:ue_sketch}): \emph{toward}, $|\Delta\phi|<60^\circ$,
which contains the leading jet; \emph{away}, $|\Delta\phi|>120^\circ$, which
contains its recoil; and \emph{transverse}, $60^\circ<|\Delta\phi|<120^\circ$,
which contains neither and is sensitive to the multiparton interactions,
the remnants, and the wide-angle part of the initial-state shower. The
observable is the density
\begin{equation}
  \rho_{\mathrm{UE}}=\frac{\sum_{\mathrm{transverse}}\pt}{\Delta\eta\,\Delta\phi}
  =\frac{\sum_{\mathrm{transverse}}\pt}{(2\times2.5)\times(2\pi/3)},
  \label{eq:uninvited_cousins:density}
\end{equation}
per unit of $\eta$--$\phi$ area, for charged particles with $|\eta|<2.5$ in
the measurements and for the final partons here. A tune~\cite{Skands:2014pea,CMS:2019csb}
is the fit of $\ptzero^{\mathrm{ref}}$, $\epsilon$, the proton profile and
the reconnection range to this density and its growth with the leading
object's \pt: more screening (larger \ptzero) means fewer scatterings and
a lower plateau, a more compact core means a larger spread between central
and peripheral collisions.

\begin{figure}[htbp]\centering
\resizebox{\textwidth}{!}{% ---------------------------------------------------------------------------
% The transverse-region definition of the underlying event (CDF, Field):
% azimuth measured from the leading object (here the b quark, phi = -2.73),
% toward |dphi| < 60 deg, transverse 60-120 deg, away > 120 deg, within
% |eta| < 2.5. Numbers of our event from the -v log summary: toward 49.47,
% transverse 21.99, away 47.67 GeV of partons not from the hard system.
% ---------------------------------------------------------------------------
\begin{tikzpicture}[x=1cm,y=1cm,
  every node/.style={font=\footnotesize},
  lab/.style={font=\scriptsize,align=center},
  head/.style={font=\small\bfseries}]
% ---------------- (a) the azimuthal regions ---------------------------------------------------
\node[head,anchor=west] at (-3.2,3.6) {(a) three regions in azimuth around the leading $b$};
\fill[red!14] (0,0) -- (60:3) arc[start angle=60,end angle=-60,radius=3] -- cycle;
\fill[teal!18] (0,0) -- (120:3) arc[start angle=120,end angle=60,radius=3] -- cycle;
\fill[teal!18] (0,0) -- (-60:3) arc[start angle=-60,end angle=-120,radius=3] -- cycle;
\fill[blue!10] (0,0) -- (120:3) arc[start angle=120,end angle=240,radius=3] -- cycle;
\draw[thick] (0,0) circle (3);
\foreach \a in {60,120,240,300} {\draw[dashed,black!60] (0,0) -- (\a:3);}
\draw[very thick,red!80!black,-{Stealth[length=2.5mm]}] (0,0) -- (0:2.6) node[right,font=\scriptsize,text=red!80!black] {leading $b$, $\pt=78\GeV$};
\node[lab,text=red!80!black] at (0:1.3) {toward\\$|\Delta\phi|<60^\circ$};
\node[lab,text=teal!80!black] at (90:1.9) {transverse\\$60^\circ<|\Delta\phi|<120^\circ$};
\node[lab,text=teal!80!black] at (-90:1.9) {transverse};
\node[lab,text=blue!70!black] at (180:1.5) {away\\$|\Delta\phi|>120^\circ$};
\node[lab,anchor=north,text=black!65] at (0,-3.3) {only partons with $|\eta|<2.5$ and not descending from the hard system count};
% ---------------- (b) the numbers of our event ------------------------------------------------
\begin{scope}[xshift=7.6cm,yshift=-3.3cm]
\node[head,anchor=west] at (-0.3,6.9) {(b) our event and the density};
\draw[-{Stealth}] (0,0) -- (6.0,0);
\draw[-{Stealth}] (0,0) -- (0,5.6) node[above,font=\scriptsize] {$\sum\pt$ of UE partons [GeV]};
\foreach \x/\l in {1.0/toward,3.0/transverse,5.0/away} {\node[font=\scriptsize,anchor=north] at (\x,-0.1) {\l};}
\foreach \y/\l in {1/10,2/20,3/30,4/40,5/50} {\draw (0,\y) -- (-0.1,\y) node[left] {\l};}
\fill[red!40] (0.5,0) rectangle (1.5,4.947);
\fill[teal!50] (2.5,0) rectangle (3.5,2.199);
\fill[blue!30] (4.5,0) rectangle (5.5,4.767);
\node[font=\scriptsize,anchor=south] at (1.0,4.95) {$49.5$};
\node[font=\scriptsize,anchor=south] at (3.0,2.2) {$22.0$ (21 partons)};
\node[font=\scriptsize,anchor=south] at (5.0,4.77) {$47.7$};
\node[font=\scriptsize,anchor=north west,align=left,text width=5.6cm] at (0.2,-0.6) {transverse density: $\dfrac{21.99\GeV}{2\times(2\times2.5)\times(\pi/3)}=\dfrac{21.99}{10.47}=2.10\GeV$ per unit $\eta$--$\phi$;\\ toward and away carry the recoil partons of the hard shower (ISR), transverse does not};
\node[font=\scriptsize,anchor=north west,align=left,text=black!65,text width=5.6cm] at (0.2,-3.0) {ensembles: $3.23\pm1.44$ (this script), $3.48\pm1.54$ (\PYTHIA) at our $b$; $2.00$ and $2.21$ with $b$ drawn};
\end{scope}
\end{tikzpicture}
}
\caption[The transverse-region definition of the underlying event]{Left:
the three azimuthal regions relative to the leading $b$ quark of our event.
Right: the summed \pt of the partons not from the hard system in each
region of our event, and the transverse density of
\cref{eq:uninvited_cousins:density}.}
\label{fig:uninvited_cousins:ue_sketch}
\end{figure}

\paragraph{The numbers.} In our event the partons that do not descend from
the hard system carry $49.5\GeV$ in the toward region, $22.0\GeV$ in the
transverse region, $47.7\GeV$ in the away region, within $|\eta|<2.5$. The
toward and away sums are large because the initial-state partons of the
hard system are counted here as ``not from the hard system'' in the
analyser's sense while the hard shower's recoil lives in those regions; the
transverse region holds 21 partons and
$\rho_{\mathrm{UE}}=21.99/10.47=2.10\GeV$. Over the ensembles the transverse
density is $3.23\pm1.44\GeV$ at our impact parameter and $2.00\pm1.59$ with
$b$ drawn, against \PYTHIA's $3.48$ and $2.21$; our event is a low
fluctuation at its $b$, and the all-azimuth density of its secondary
partons, $3.79\GeV$, is closer to the ensemble mean. The measured
plateau at $13\TeV$ for charged particles above $0.5\GeV$ is about
$1\GeV$ per unit area in the transverse region~\cite{ATLAS:2017blj}, which
after the thresholds and the neutral fraction is consistent with the
parton-level numbers here; the comparison that matters is the one the tune
makes, with \PYTHIA at hadron level, and CP5 reproduces the data to within
its uncertainties~\cite{CMS:2019csb}. Panel~(c) of
\cref{fig:uninvited_cousins:ue_plot} is the tune knob in action:
$\langle n\rangle$ against \ptzero from the script's \code{SigmaInt},
falling from $25$ at $\ptzero=1\GeV$ to $3$ at $2.5\GeV$.

\begin{figure}[htbp]\centering
\includegraphics[page=7,width=\textwidth]{MPI/Standalone/theory_multiparton_interactions.pdf}
\caption[Underlying-event densities and the tune knob]{Page~7 of
\file{theory_multiparton_interactions.pdf}. (a) The transverse and
all-azimuth densities of the partons not from the hard system for the two
standalone and the two \PYTHIA ensembles, with our event as the star.
(b) The event-by-event distribution of the transverse density. (c) The mean
number of scatterings against the regulator \ptzero, and $f(b)\langle n\rangle$
at our impact parameter.}
\label{fig:uninvited_cousins:ue_plot}
\end{figure}

\section{Validation from official packages}
\label{sec:uninvited_cousins:validation}

Two official sources again: \PYTHIA driven from a script so that it dresses
exactly our hard event, and the CMS generator job of \cref{ch:ancestral_home},
which had multiparton interactions, remnants and colour reconnection on
from the start. As before, ensembles are compared on averages and the
single \CMSSW event on structure.

\subsection{The same stage with \PYTHIA: the validation script}
\label{sec:uninvited_cousins:pythia}

\file{validate_mpi_with_pythia.py} writes the hard event of
\cref{ch:ancestral_home} into a Les Houches file $N$ times with
\code{SCALUP} $=113.32\GeV$ and lets \PYTHIA~8.315, in the environment of
\file{example/ME/Standalone/}, run its parton level with the CP5 settings of
the \CMSSW configuration: showers, multiparton interactions, remnants and
colour reconnection on, hadronisation and QED off. It reads \PYTHIA's
initialisation numbers, assigns every final parton to its system by walking
the record, and measures the same quantities as the standalone
\code{--repeat}. Two things in it are worth knowing before you write your
own. \code{Tune:pp = 14} resets \code{PDF:pSet}, so the tune must be set
\emph{before} the PDF, as the \CMSSW configuration does; the chapter-2
validation script had the order reversed and its \PYTHIA numbers moved by
one to two percent when it was corrected. And \PYTHIA prints its
initialisation through C++ streams, invisible to Python's \code{sys.stdout},
so the file descriptor is redirected while \code{init()} runs. The impact
parameter cannot be fixed from the Python binding, so the fixed-$b$ sample
keeps the events with $|b/\bmean-0.3136|<0.03$, $537$ of $8000$:
\begin{minted}[linenos=false,frame=single,framesep=1.5mm,fontsize=\small]{bash}
cd example/MPI/Standalone
PY=../../ME/Standalone/.venv/bin/python
$PY validate_mpi_with_pythia.py -n 2000 --standalone-stats standalone_mpi_stats_sampled.stats.json
$PY validate_mpi_with_pythia.py -n 8000 --impact 0.3136 --window 0.03
$PY validate_mpi_with_pythia.py -n 500 --no-cr
\end{minted}
\PYTHIA dresses two thousand copies in about a minute and prints
$\ptzero=1.44\GeV$, $\sigint=496.6\,\mathrm{mb}$ and $\signd=55.42\,\mathrm{mb}$
at initialisation, against the script's $497.3$; the \CMSSW log's $500.67$
differs from both by the version of the integration grid. The columns of
\cref{tab:uninvited_cousins:numbers} are the result. Read them as before.
The geometry agrees: $\langle b\rangle$, $\langle f\rangle$ and $f(0.3136)$
within $2\%$. The remnant count, the hardest secondary scattering and the
underlying-event density agree at the $5$--$10\%$ level, and the extra \pt
inside the cones follows the density. The one real disagreement is the
number of scatterings, $20\%$ low in the standalone chain at both
geometries, and its cause is known: \PYTHIA interleaves the multiparton
interactions with the initial- and final-state showers of \emph{all}
systems in one common \pt sequence, so that the depletion a scattering sees
is only that of the scatterings and emissions harder than itself, while
the script showers each system to its cut-off before drawing the next,
softer scattering, which makes the proton look emptier than it should.
Without any showering the two chains agree to $4\%$ ($7.2$ against $7.5$
scatterings per unit of $f$). The remaining difference, a $10\%$ stronger
soft radiation of the chapter-2 shower, was already seen in
\cref{sec:mother_gives_birth:pythia}. The cone \emph{mass} is not compared:
the standalone ensembles keep the seed-1 shower, whose family mass of
$13.5\GeV$ is a high draw, while \PYTHIA reshowers to $8\GeV$ on average.
\PYTHIA's first dressed event is in
\file{validate_mpi_with_pythia_event1.txt} in the standalone format.

\subsection{How CMS does it in practice: the rest of the protons inside the GEN step}
\label{sec:uninvited_cousins:cmssw}

Nothing new has to be run: the \code{GEN} step of
\cref{sec:ancestral_home:cmssw} already produced this stage, and
\code{JitJetSW}'s analyser drew it as its third figure. From the setup of
\cref{sec:mother_gives_birth:cmssw},
\begin{minted}[linenos=false,frame=single,framesep=1.5mm,fontsize=\small]{bash}
cd CMSSW_15_0_5/src/JitJetSW/Journey/test
./runCMSDrivers_mc2024_BBbar_GEN.sh          # GEN step + genPartAnalyzer, one PDF per stage
\end{minted}
produces \file{genPartAnalyzer_ME_PS_MPI_out.pdf}, whose third page is
\cref{fig:uninvited_cousins:cmssw}; the copy in \file{example/MPI/CMSSW/}
is the one used here, and the log is the one of \cref{ch:mother_gives_birth}.

\subsubsection{Reading the log}

Four lines of \file{step1_GEN_PS_cfg.log} are this chapter. At
initialisation,
\begin{minted}[linenos=false,frame=single,framesep=1.5mm,fontsize=\small]{text}
 |    pT0 =  1.44 gives sigmaInteraction =   500.67 mb: accepted    |
\end{minted}
is \cref{eq:uninvited_cousins:regularised,eq:uninvited_cousins:dsigma}: \PYTHIA
integrates the regularised cross section at the start of every job and
accepts the regulator if the integral is large enough for the veto
algorithm's overestimate; the warning that follows,
``\code{MultipartonInteractions::init: maximum increased by factor 1.090}'',
is the overestimate of \cref{eq:uninvited_cousins:trial} being tightened
by a trial scan. The settings table below it lists \code{pT0Ref = 1.41},
\code{ecmPow = 0.03344}, \code{coreFraction = 0.63}, \code{coreRadius = 0.7634},
\code{bProfile = 2} and \code{ColourReconnection:range = 5.176}, the CP5
numbers used throughout. In the Info Listing of the first event, after the
hard-process invariants,
\begin{minted}[linenos=false,frame=single,framesep=1.5mm,fontsize=\small]{text}
 Impact parameter b =  3.136e-01 gives enhancement factor =  3.709e+00.
 Max pT scale for MPI =  1.133e+02, ISR =  1.133e+02, FSR =  1.133e+02.
 Number of MPI =    32, ISR =    28, FSRproc =    64, FSRreson =     0.
\end{minted}
The first line is \cref{eq:uninvited_cousins:fb} for this event, the value
the book's record was fixed to; the second is the common starting scale of
the interleaved evolution; the third counts the hard process among the
``MPI'', so the event has 31 secondary scatterings, and its 28 initial-state
and 64 final-state branchings belong to all 32 systems together, which is
why they could not be compared with \cref{tab:mother_gives_birth:numbers}
at the previous stop. In the drawer's table the new status codes of
\cref{sec:uninvited_cousins:record_format} appear in the same roles: 62
rows of status~31 and 62 of status~33, the incoming and outgoing partons of
the 31 secondary scatterings; 64 rows of status~61, the initiators of all
32 systems after the remnant kinematics; 156 of status~62; 13 of status~63,
the remnants; and then 120 of status~71 and 45 of status~73, the copies
that \PYTHIA makes when it collects each colour-singlet string and joins
very soft partons for hadronisation, which the standalone record does not
have. Follow the mother once more: row~515 is her status-62 copy at
$91.744\GeV$, $0.15\GeV$ above the $91.59$ of her last shower copy (row 182,
status~52), the same primordial-$k_{\mathrm{T}}$ boost that moved ours from
$78.12$ to $78.29$; row~824, status~71, hands her to hadronisation at the
same momentum.

\subsubsection{The figure}

\Cref{fig:uninvited_cousins:cmssw} is the analyser's stage-3 picture: 169
final partons, 61 from the secondary scatterings by the analyser's rule, 13
remnants of which most are on the border at $|\eta|>6$, and the $b$ and
$\bar b$ where the previous stage left them. Set it next to
\cref{fig:uninvited_cousins:journey}. The two are different dressings of
different showers of the same hard event, and the comparison to make is
\cref{tab:uninvited_cousins:numbers}; what the pictures share is the
structure, cyan squares spread uniformly in $\eta$ and $\phi$ with no
preferred direction, far more of them at $|\eta|>3$ than a central detector
will ever see, and a handful of them inside the cones that the later
chapters will draw around the two $b$ quarks.

\begin{figure}[htbp]\centering
\includegraphics[page=3,width=\textwidth]{MPI/CMSSW/genPartAnalyzer_ME_PS_MPI_out.pdf}
\caption[The event after MPI and beam remnants in \CMSSW]{Stage~3 of our
event in \CMSSW, drawn by the \code{genPartAnalyzer} of \code{JitJetSW}:
the parton level after multiparton interactions, beam remnants and colour
reconnection. Red and blue are the status-62 copies of the $b$ and $\bar b$,
orange and green triangles the initial- and final-state emissions of all
systems, cyan squares the outgoing partons of the 31 secondary scatterings
and their copies, grey diamonds the 13 beam remnants; partons beyond
$|\eta|=6$ are drawn on the border. The same stage of the standalone
record is \cref{fig:uninvited_cousins:journey}.}
\label{fig:uninvited_cousins:cmssw}
\end{figure}

\section{Our jet at this stage: UE energy inside \texorpdfstring{$R=0.4$}{R = 0.4} and \texorpdfstring{$R=0.8$}{R = 0.8}}
\label{sec:uninvited_cousins:our_jet_at_this_stage_ue_energy_inside_r}

Here is the family among the cousins. The record,
\file{example/MPI/event_MPI.lhe}, has 963 rows. The first 158, rows 0 to
157, are the showered event of \cref{ch:mother_gives_birth}, every one of
its final partons now with a negative status because the stage copied
them; rows 158 to 161 are the first secondary scattering:
\filelistinglines[\widelisting]{text}{example/MPI/event_MPI.lhe}{162}{165}{the first secondary scattering, $gg\to gg$ at $\pt=11.45\GeV$: two incoming gluons (status $-31$) and two outgoing (status $-33$), with the colour tags of the log}
whose outgoing gluons have $p_x=\pm3.302$, $p_y=\pm10.967\GeV$, that is
$\pt=11.453\GeV$, and whose incoming gluons carry $p_z=3.317$ and
$-103.979\GeV$, that is $x_1=4.88\times10^{-4}$ and $x_2=1.53\times10^{-2}$
of $6800\GeV$. The next hundred rows are its shower and the 23 further
scatterings with theirs. From row 702 the stage writes the boosted copies:
702 and 703 are the hard system's initiators after their kicks
(status $-61$, $p_x,p_y=-1.767,+1.652$ and $-1.113,-1.838\GeV$), and far
below, among the status-62 copies,
\filelistinglines[\widelisting]{text}{example/MPI/event_MPI.lhe}{723}{723}{the mother's final copy (row 719, status 62) after the primordial-$k_{\mathrm{T}}$ boost}
is the mother at $78.29\GeV$, mother row 124, her last shower copy. The
sixteen remnants close the record,
\filelistinglines[\widelisting]{text}{example/MPI/event_MPI.lhe}{951}{953}{the first three remnants of proton A (status 63): the $ud_0$ diquark and two companion $u$ quarks}
\filelistinglines[\widelisting]{text}{example/MPI/event_MPI.lhe}{964}{967}{the last two remnants of proton B and the momentum sum over the final state}
and the momentum sum is that of the two protons.

\begin{table}[htbp]\centering\small
\caption[Kinematics of the two $b$ quarks and their cones after the MPI
stage]{Kinematics of the two $b$ quarks and of the partons within $\dR<0.4$
and $0.8$ of each after the MPI stage, from \file{event_MPI.lhe} as printed
by \file{standalone_multiparton_interactions.py}. ``Hard system'' means the
partons that descend from the hard process and its shower; ``UE'' the
partons of the secondary systems and the remnants.}
\label{tab:uninvited_cousins:kinematics}
\footnotesize
\begin{tabular}{llrrrr}\toprule
 & & $\pt$ [GeV] & $m$ [GeV] & $n$ & of which UE \\ \midrule
$b$ & after the shower (chapter 2, row 124) & $78.12$ & $4.80$ & $1$ & \\
    & after the kick (row 719, status 62) & $78.29$ & $4.80$ & $1$ & \\
    & hard system within $\dR<0.4$ & $100.67$ & $13.52$ & $3$ & \\
    & all partons within $\dR<0.4$ & $100.67$ & $13.52$ & $3$ & $0$ \\
    & hard system within $\dR<0.8$ & $100.67$ & $13.52$ & $3$ & \\
    & all partons within $\dR<0.8$ & $104.09$ & $20.17$ & $6$ & $3$ ($3.42\GeV$) \\ \midrule
$\bar b$ & after the kick (status 62) & $72.50$ & $4.80$ & $1$ & \\
    & all partons within $\dR<0.4$ & $93.88$ & $7.75$ & $4$ & $1$ ($2.66\GeV$) \\
    & all partons within $\dR<0.8$ & $116.07$ & $41.61$ & & $3$ \\ \midrule
all final partons & & $0.00$ & $13600$ & $211$ & $171$ \\
\bottomrule
\end{tabular}
\end{table}

\Cref{tab:uninvited_cousins:kinematics} and \cref{fig:uninvited_cousins:journey}
say what the record means for the ledger. Within $\dR<0.4$ of the mother
nothing arrived: the three partons of her family are still her two
children and herself, and the only change is the $0.22\GeV$ of the
primordial-$k_{\mathrm{T}}$ boost. That is luck, and
\cref{fig:uninvited_cousins:our_jet_plot} says how much: at this impact
parameter a cone of $R=0.4$ receives $1.6\pm2.2\GeV$ on average, one parton
in most events and none in a third of them; a uniform density of
$3.79\GeV$ per unit area, the all-azimuth density of our event, would put
$\rho\pi R^2=1.9\GeV$ into it. Within $\dR<0.8$ three cousins did arrive,
with $3.42\GeV$ of \pt between them, and they raise the family mass from
$13.5$ to $20.2\GeV$: the mass is far more sensitive to a soft parton at a
large angle than the \pt is, because the mass grows with the product of
the energy and the angle. The $\bar b$ at $\eta=-3.4$ received one cousin
of $2.66\GeV$ already at $R=0.4$. This is the general lesson of the stage,
and it will return in the detector chapters with the pileup in place of the
cousins: a cone of radius $R$ collects an energy that grows as $R^2$ from
everything that is uniform in the event, while the family's own energy
saturates once $R$ exceeds the shower's width
(\cref{fig:mother_gives_birth:runaways_plot}c). The choice of $R=0.4$ for CMS
jets is the compromise between the two curves.

\begin{figure}[htbp]\centering
\includegraphics[page=8,width=\textwidth]{MPI/Standalone/theory_multiparton_interactions.pdf}
\caption[The underlying-event energy inside the $b$ cone]{Page~8 of
\file{theory_multiparton_interactions.pdf}. (a) The scalar \pt of the
underlying-event partons inside $\dR<0.4$ and $0.8$ around the $b$ for the
four ensembles, our event as the star, and $\rho\pi R^2$ with the
all-azimuth density of our event as the horizontal lines. (b) The
event-by-event distribution: sometimes nothing, sometimes a lot. (c) The
ledger of our event against the cone radius: the vector-summed \pt of the
hard-system partons and of all partons inside $R$, and the mass of each.}
\label{fig:uninvited_cousins:our_jet_plot}
\end{figure}

\begin{figure}[htbp]\centering
\includegraphics[width=\textwidth]{MPI/journey_MPI.pdf}
\caption[The event after the MPI stage in the $(\eta,\phi)$ plane]{The
parton level of our event after the standalone MPI stage, drawn by
\file{example/MPI/plot_journey_MPI.py} from \file{example/MPI/event_MPI.lhe}
in the layout of \cref{fig:uninvited_cousins:cmssw}. Red and blue are the
status-62 copies of the $b$ and $\bar b$, magenta their vector sum; orange
and green triangles the initial- and final-state emissions of all systems,
cyan squares the outgoing partons of the 24 secondary scatterings and their
copies, grey diamonds the 16 beam remnants, most of them on the border. The
header lists the number and the net and scalar \pt of each group; the
vector sum of all 211 partons vanishes.}
\label{fig:uninvited_cousins:journey}
\end{figure}

The ledger gets its first row that adds instead of subtracting. The $R=0.4$
row is the one the later stages are compared with; the $R=0.8$ row is
listed because the large-radius jets of \cref{ch:three_families} will need
it.

\begin{ledger}{Cousins Nobody Invited}
  \ledgerrow{$b$ family after the shower ($\dR<0.4$)}{100.45}{13.52}{3}{--}
  \ledgerrow{$b$ family after this stage ($\dR<0.4$)}{100.67}{13.52}{3}{$+0.22$}
  \ledgerrow{\quad of which the $b$ quark itself}{78.29}{4.80}{1}{--}
  \ledgerrow{\quad the same within $\dR<0.8$ (3 UE partons)}{104.09}{20.17}{6}{$+3.42$}
\end{ledger}

\section{What this stage is, and what it is not}
\label{sec:uninvited_cousins:caveats}

The agreement in \cref{tab:uninvited_cousins:numbers} was obtained by
porting the generator's construction, not by tuning to its output, and the
disagreements are understood. Each omission is a caveat to carry:
\begin{itemize}
  \item \textbf{Sequential, not interleaved.} \PYTHIA runs multiparton
    interactions, ISR and FSR of all systems in one \pt-ordered sequence
    \cite{Sjostrand:2004ef}; the script showers each system before drawing
    the next scattering. The chain is $20\%$ short because of it, and the
    recoils between systems are not shared as \PYTHIA shares them.
  \item \textbf{At most one valence quark leaves a proton.} A second one is
    reclassified as sea, so the junction topologies of \PYTHIA's remnants,
    a baryon number carried by three strings meeting at a point, never
    occur. They are rare at our $x$ and matter for baryon production in
    the forward region.
  \item \textbf{No rescattering, no diffraction, no photons.} An outgoing
    parton of one scattering can scatter again in \PYTHIA's extended
    model~\cite{Corke:2009tk} (off by default, and in CP5); the elastic and
    diffractive $45\,\mathrm{mb}$ of the total cross
    section~\cite{Navin:2010kk} are subtracted from \signd and never
    simulated here; photon-induced
    processes are absent.
  \item \textbf{The MPI-based colour reconnection only.} \PYTHIA's newer
    QCD-based model~\cite{Christiansen:2015yqa} and the gluon-move model
    are not ported; CP5 uses the MPI-based one with $R=5.176$, so the
    comparison is like with like, but the open problems below are about
    exactly this choice.
  \item \textbf{A fixed proton size.} The double Gaussian of CP5 has no
    $x$ dependence; \PYTHIA offers an $x$-dependent proton
    width~\cite{Corke:2011yy} that CMS's other tunes do not use either.
  \item \textbf{Parton level only.} Colour reconnection moves no momentum
    here; its effect, and the remnants' and the cousins' hadrons, appear in
    \cref{ch:children_getting_married}. The ``partons'' below $1\GeV$ are string
    pieces, not objects.
  \item \textbf{The impact parameter is fixed by hand} in the book's record.
    This is a choice for the sake of comparison with the \CMSSW event; the
    ensembles with $b$ drawn are the physics.
  \item \textbf{One dressing is one sample.} The $0.00\GeV$ inside the
    $R=0.4$ cone is a draw from a distribution whose mean is $1.6\GeV$; the
    ledger rows are those of \emph{this} event.
\end{itemize}

\section{Further reading}
\label{sec:uninvited_cousins:further_reading}
The multiple-interaction model with its impact-parameter dependence is
\cite{Sjostrand:1987su}; the depleted proton, the beam remnants and the
original colour-reconnection rule are \cite{Sjostrand:2004pf}; the
interleaving with the showers is \cite{Sjostrand:2004ef}; Sj\"ostrand's own
history of the modelling, with the reasons for every choice, is
\cite{Sjostrand:2017cdm}, and the \PYTHIA~8.3 guide~\cite{Bierlich:2022pfr}
documents every setting named here. The Herwig model, built on the same
eikonal idea with a different proton, is \cite{Bahr:2008dy}. The CDF
definition of the transverse region is \cite{CDF:2001onq,Field:2002vt}, the
$13\TeV$ measurements that constrain the tunes are
\cite{ATLAS:2017blj,CMS:2017ngy}, and the CMS tunes themselves, CP5 among
them, are \cite{CMS:2019csb}, with the colour-reconnection variants in
\cite{CMS:2022awf}; the Monash tune~\cite{Skands:2014pea} is the \PYTHIA
default they start from. Colour reconnection and precision measurements
are discussed in \cite{Sjostrand:2013cya,Argyropoulos:2014zoa}, the models
beyond leading colour in \cite{Christiansen:2015yqa,Gieseke:2012ft}. The
general review remains \cite{Buckley:2011ms}, and the jet-area method that
will subtract the uniform part of the underlying event together with the
pileup is \cite{Cacciari:2008gn}.

\begin{pitfalls}
\begin{itemize}[nosep]
  \item \textbf{Underlying event is not pileup.} The cousins of this
    chapter come from the \emph{same} proton--proton collision as the
    mother; pileup (\cref{ch:crowd_on_the_road}) is other collisions in the
    same bunch crossing. The first is physics of the event and is in every
    generator sample; the second is a detector condition and is added
    afterwards. Both deposit energy uniformly, and the subtraction methods
    of \cref{ch:paying_for_uninvited} remove both, but a tune is about the
    first only.
  \item \textbf{``Number of MPI'' includes the hard process.} \PYTHIA's
    counter, the \CMSSW log's $32$, counts the hard scattering as the first
    interaction. The event has 31 secondary scatterings.
  \item \textbf{Set the tune before the PDF.} \code{Tune:pp} resets
    \code{PDF:pSet} to the tune's own PDF. Every CMS configuration has the
    order right; a hand-written \PYTHIA card may not, and the symptom is a
    few percent everywhere and no error message.
  \item \textbf{The enhancement factor is not the overlap divided by its
    mean.} $f(b)$ is \cref{eq:uninvited_cousins:fb}; a plausible-looking
    alternative differs by a few percent at small $b$ and would make the
    \CMSSW number $3.709$ unreproducible.
  \item \textbf{Two different partons, one final state, counted twice.}
    Integrating $y_3$ and $y_4$ over their full ranges reaches $qg\to qg$
    both as $t$- and as $u$-channel; halve those channels or \sigint comes
    out $43\%$ high.
  \item \textbf{Freeze, do not extrapolate, the PDFs below the grid.}
    \PYTHIA's readers freeze the densities below $Q=1.65\GeV$; LHAPDF's own
    extrapolation gives $352$ instead of $500\,\mathrm{mb}$, and two thirds
    of \sigint is below $1.5\GeV$.
  \item \textbf{The impact parameter cannot be fixed from Python.} The
    \code{pythia8mc} wheel offers no \code{UserHooks} override; select
    events in a window of $b$ instead, as the validation script does, and
    accept the statistics.
  \item \textbf{Counting ``MPI partons'' depends on the rule.} The record
    has 155 partons from the secondary systems; the analyser's rule, which
    assigns emissions to ISR and FSR whatever system made them, counts 45.
    Say which rule you use before comparing with a \CMSSW figure.
  \item \textbf{Nothing in the cone is luck, not physics.} At this impact
    parameter the $R=0.4$ cone receives $1.6\pm2.2\GeV$; our $0.0$ is one
    draw. Do not conclude from one event that the underlying event is
    negligible at $R=0.4$, and do not conclude from the mass jump at
    $R=0.8$ that it is large: both are the tails.
  \item \textbf{The remnants are at $|\eta|>6$, mostly.} Fifteen of the
    sixteen remnant partons are beyond any detector; the one at
    $\eta=0.35$ carries $0.67\GeV$. The remnants matter for the energy flow
    and the forward detectors, not for the barrel jet.
  \item \textbf{Colour reconnection is invisible at parton level.} Every
    momentum in the record is the same with \code{--no-cr}; only the colour
    tags differ. Its effects are in the hadrons, and a parton-level
    comparison cannot validate a reconnection model.
  \item \textbf{Which Python?} The standalone, plotting and animation
    scripts run with a \code{python3} that has NumPy, SciPy and matplotlib;
    the validation script needs the virtual environment of
    \file{example/ME/Standalone/}. The scripts say what they loaded on their
    first line.
\end{itemize}
\end{pitfalls}

\begin{exercises}
\item \textbf{Run it.} Run \file{standalone_multiparton_interactions.py}
  with \code{--impact 0.3136} and confirm $24$ secondary scatterings, a
  transverse density of $2.10\GeV$ and $0.00\GeV$ of underlying event
  inside $\dR<0.4$ of the $b$. Run it without \code{--impact} for seeds 1
  to 10 and tabulate $b/\bmean$, $f(b)$ and the number of scatterings;
  check that the ratio of the last two is roughly constant.
\item \textbf{The divergence.} From the table of
  \cref{sec:uninvited_cousins:cross_section}, estimate \sigint if the
  regulator were removed and the integral cut at $\pt^{\min}=0.2$, $0.5$
  and $1\GeV$ instead (the bare rate is the regularised one divided by the
  third column). Compare with \signd and explain why none of the three
  cut-offs is an acceptable model.
\item \textbf{The energy dependence.} Compute \ptzero at $7$, $13$, $13.6$
  and $14\TeV$ from \cref{eq:uninvited_cousins:regularised} and, with
  \code{SigmaInt} from the script, \sigint and $\langle n\rangle$ at each
  (change \code{ECM} and the \code{--sigma-nd} input; the total cross
  sections are in the \PYTHIA guide). By how much does the average number
  of scatterings change between $13.6$ and $14\TeV$?
\item \textbf{The overlap by hand.} Write $O(b)$ of
  \cref{eq:uninvited_cousins:overlap} for the double Gaussian with
  $\beta=0.63$ and $a_2=0.7634\,a_1$, normalise it, and reproduce $f(0)=4.52$
  and $f(0.3136\bmean)=3.63$ with \cref{eq:uninvited_cousins:fb} and a
  one-dimensional numerical integral in $b$. Then compute $\langle f\rangle$
  with the minimum-bias and with the hard-process weight and confirm
  $0.97$ and $2.22$.
\item \textbf{Invert the chain.} Derive \cref{eq:uninvited_cousins:trial}
  from \cref{eq:uninvited_cousins:chain} with the overestimate
  $g(\pt^2)=C/(\pt^2+\ptzero^2)^2$, $C=9752\,\mathrm{mb}\,\GeV^2$ from the
  log. Starting at $\pt^{\mathrm{prev}}=113.3\GeV$ with $f=3.627$, compute
  the median trial scale and compare with the $11.45\GeV$ at which the
  first scattering was accepted; then explain, with the veto weight
  $0.174$, why $303$ trials were needed for $24$ acceptances.
\item \textbf{The depleted proton.} With \code{Beam} from the script,
  reproduce the modified densities of proton B at the first rung,
  $x/X=0.0289$: $xg=4.26$, $xu=0.487$, $xd=0.537$, and show that the valence
  $u$ density has been halved and the gluon and sea multiplied by $1.164$
  relative to the unmodified proton at the same $x'$. What would the
  factor be if both $u$ quarks had left?
\item \textbf{Companions.} Plot the companion distribution
  $f_c(x_c;x_s)$ for $x_s=0.02$ and $0.2$ and explain why a companion of a
  large-$x$ sea quark is itself likely to carry little momentum. Find the
  six companions of proton A in the remnant rows of \file{event_MPI.lhe} and
  compare their $x$ values with the curve.
\item \textbf{The kick.} From rows 702 and 703 of \file{event_MPI.lhe}
  verify the two kicks of the hard system and compute, by boosting the
  $b$ of row 124 to the frame of the kicked initiators, its new \pt;
  confirm $78.29\GeV$. Then rerun with \code{--no-kt} and check that the
  family \pt returns to $100.45\GeV$.
\item \textbf{Light-cone closure.} Sum $p^+=E+p_z$ over the status-62
  partons of the systems that took their initiator from proton A and over
  the remnants of A (rows 947 to 953), and confirm that the total is
  $13600\GeV$; repeat for $p^-$ and proton B. Which remnant carries the
  largest share?
\item \textbf{Reconnection by hand.} For rung 2 ($\pt=7.679\GeV$) and rung 1
  ($11.453\GeV$) compute \cref{eq:uninvited_cousins:cr_prob} with
  $R=5.176$ and confirm $0.486$ and $0.298$. Then take the $7.34\GeV$ gluon
  of rung 2 (row 778) and the dipole it was inserted into (rows 702 and
  738) and verify the invariant $441.8\GeV^2$ of
  \cref{eq:uninvited_cousins:insertion}. Run with \code{--cr-range 1.8},
  the Monash value, and tabulate how many systems merge and the final
  $\lambda$.
\item \textbf{The string length.} Compute $\lambda$ of
  \cref{eq:uninvited_cousins:lambda} for the hard system alone from the
  rows of \file{event_PS.lhe} and its colour tags, and compare with the
  per-system contributions the script prints with \code{-vv}. Why does
  merging a soft system into a hard one \emph{reduce} the total?
\item \textbf{The transverse region.} From the status-62 and 63 rows of
  \file{event_MPI.lhe}, select the partons that do not descend from the
  hard process (walk the mothers), restrict to $|\eta|<2.5$, and reproduce
  the three sums $49.5$, $22.0$ and $47.7\GeV$ and the density $2.10\GeV$.
  Repeat with \emph{all} partons: by how much do the toward and away
  regions grow, and why does the transverse region barely change?
\item \textbf{The cone against the radius.} Reproduce panel~(c) of
  \cref{fig:uninvited_cousins:our_jet_plot} from the record with the reader
  of \file{plot_journey_MPI.py}: the vector-summed \pt and the mass of the
  hard-system partons and of all partons within $R$ of the $b$, for $R$
  from $0$ to $2$. At which $R$ does the first cousin enter, and what is it?
  Then use \code{--repeat 100} with \code{--radius 0.2}, \code{0.4},
  \code{0.8} and \code{1.0} and plot the mean underlying-event \pt against
  $R^2$.
\item \textbf{Read the \CMSSW record.} In \file{step1_GEN_PS_cfg.log} list
  the 31 status-33 pairs with their \pt, reproduce the ladder of
  \cref{fig:uninvited_cousins:chain_plot}b, and find the 13 status-63
  remnants with their $\eta$. Which flavours did the remnants of each proton
  get, and what does that tell you about which valence quarks left?
\item \textbf{\PYTHIA's first event.} Run \file{plot_journey_MPI.py} on
  \file{validate_mpi_with_pythia_event1.txt} and compare the header with
  \cref{fig:uninvited_cousins:journey}. Then run the validation script
  with \code{--no-cr} for 500 events and verify that no momentum-based
  quantity changes beyond its statistical error.
\item \textbf{Film a peripheral collision.} Run the animation with
  \code{--impact 1.2 --scene chain} and, before watching, predict from
  \cref{fig:uninvited_cousins:overlap_sketch}b how many rungs the ladder
  will have. Then \code{--scene strings --no-cr} and describe what the
  right-hand bar does.
\end{exercises}

\subsection*{Open problems in the context of Phase-2 LHC}
\label{sec:uninvited_cousins:open_problems}

The underlying event is the part of the simulation that is tuned rather
than derived, and the High-Luminosity LHC will stress it from two sides:
its measurements of jets and of the top quark will be precise enough to
feel the tune's assumptions~\cite{Azzi:2019yne}, and its pileup of $140$ to
$200$ collisions per crossing will bury the underlying event under a
hundred times its own energy. Each of the following can be started with the
scripts of this chapter.
\begin{enumerate}
  \item \textbf{Colour reconnection and the top-quark mass.} The
    reconnection model shifts the reconstructed top mass by up to
    $0.5\GeV$ between \PYTHIA's three models~\cite{Argyropoulos:2014zoa},
    and the Run-2 measurements are already at $0.4\GeV$ total
    uncertainty~\cite{CMS:2023ebf,ATLAS:2024dxp}, with the reconnection
    among the largest modelling terms. CMS has tuned all three models to
    underlying-event data~\cite{CMS:2022awf} and found that the data do not
    choose between them; the jet-pull colour-flow observables of
    \cite{ATLAS:2018olo} do discriminate but with limited power. Which
    observable of a $3000\,\mathrm{fb}^{-1}$ $\ttbar$ sample would fix the
    model? Start by measuring, in the standalone script, how the string
    through the $b$ (\cref{fig:uninvited_cousins:cr_plot}c) changes with
    \code{--cr-range}, and then ask which hadron-level quantity of
    \cref{ch:children_getting_married} follows it.
  \item \textbf{Beyond leading colour.} The MPI-based rule of
    \cref{eq:uninvited_cousins:cr_prob} knows nothing about SU(3); the
    QCD-based model~\cite{Christiansen:2015yqa} forms junctions when three
    colour lines meet and predicts enhanced baryon production that ALICE
    and LHCb have seen in part, and Herwig's model~\cite{Gieseke:2012ft}
    minimises a different length. Port the QCD-based rule into
    \code{reconnect} (it needs the colour tags the record already has) and
    measure how the string length and the number of junction strings
    change at our impact parameter. What would distinguish the two rules in
    a Phase-2 measurement of baryon-to-meson ratios inside jets?
  \item \textbf{Collectivity in small systems.} The near-side ridge in
    high-multiplicity proton--proton collisions~\cite{CMS:2010ifv} and the
    strangeness enhancement with multiplicity have no explanation in the
    independent-string picture of this chapter; rope hadronisation and
    string shoving~\cite{Bierlich:2014xba,Bierlich:2017vhg} add
    interactions between overlapping strings and reproduce part of both.
    Our event has $51$ strings in a transverse area of a few fm$^2$. Using
    the strings the script prints, estimate how many overlap in the
    transverse plane at the $b$ quark's position; the question is whether
    a $b$ jet in a central collision hadronises in the same vacuum as one
    in a peripheral collision, and HL-LHC $b$-jet samples will be large
    enough to split by underlying-event activity.
  \item \textbf{Double parton scattering as a signal.} Two hard
    scatterings in one collision are the tail of the chain of
    \cref{sec:uninvited_cousins:chain}; the effective cross section
    $\sigma_{\mathrm{eff}}\approx12$--$21\,\mathrm{mb}$ measured in $W$+2~jets
    \cite{CMS:2013huw} and same-sign $WW$~\cite{CMS:2022pio} is, in the
    model, set by the overlap of \cref{eq:uninvited_cousins:overlap}, and
    the theory beyond the model, with correlated double densities, is
    reviewed in \cite{Diehl:2017wew,Bansal:2014paa}. Compute
    $\sigma_{\mathrm{eff}}=\left[\int\mathrm{d}^2b\,O(b)^2/(\int O)^2\right]^{-1}$
    for the CP5 profile with the script's \code{Overlap} and compare with
    the measurements; then ask which double-parton correlations an HL-LHC
    measurement of $\bbbar+$jets, our own process plus a cousin, could
    constrain.
  \item \textbf{A tune for $14\TeV$ and the energy dependence of
    \ptzero.} The $\epsilon=0.03344$ of CP5 was fitted to the growth of the
    underlying event between $1.96$ and $13\TeV$~\cite{CMS:2019csb}; the
    $x$-dependent proton size of \cite{Corke:2011yy} predicts the growth
    instead of fitting it. Using Exercise~3, estimate how much
    $\langle n\rangle$ and the transverse density change between $13.6$ and
    $14\TeV$ for the two models, and compare with the precision of the
    first Run-4 underlying-event measurements. Where the two models differ
    by more than that, the tune is a prediction.
  \item \textbf{Underlying event under 200 pileup.} At the HL-LHC the
    uniform energy density in a jet cone will be dominated by pileup, and
    the jet-area subtraction~\cite{Cacciari:2008gn} removes the sum of
    both. The underlying event is then absorbed into the pileup offset of
    \cref{ch:paying_for_uninvited}, with the \emph{wrong} tune
    dependence: the offset is derived from zero-bias data and the
    underlying event from the generator. With \code{--repeat} at a few
    impact parameters, quantify how much of the $1.6\GeV$ in the $R=0.4$
    cone is correlated with the hard scattering's $b$ and therefore not
    uniform; this is the part no area subtraction can remove, and its
    size at $\pt=100\GeV$ decides whether a Phase-2 jet energy scale needs
    a separate underlying-event term.
  \item \textbf{Interleaving and the budget of a $14\TeV$ event.} The
    standalone chain is $20\%$ short because it does not interleave the
    scatterings with the showers (\cref{sec:uninvited_cousins:pythia}), and
    the full interleaved evolution of \PYTHIA~\cite{Sjostrand:2004ef} is the
    single most expensive part of the parton level, which the HSF generator
    working group identifies as a cost driver for Phase-2
    production~\cite{HSFPhysicsEventGeneratorWG:2020gxw}. Implement the
    interleaving in the script by merging \code{next\_pt} with the shower's
    competition of \cref{sec:mother_gives_birth:step_by_step}, confirm that
    the chain length moves towards \PYTHIA's, and measure the cost. The
    question for Phase-2 is whether an interleaved but \emph{truncated}
    evolution, with the softest systems treated collectively, can
    reproduce the underlying-event observables at a fraction of the time.
\end{enumerate}
Two of these, the colour reconnection and the hadrons of the cousins, are
taken up at the next stop, where the partons are finally dressed.
